% concatenated from hardyInequality-Bouthat.tex
\documentclass[oneside,10pt]{article}          % please do not change
\usepackage[b5paper]{geometry}	    % your paper can be easily printed on a4 or letter paper with enlargenment      
                                    % comment if you have problem with print     
\usepackage{amsfonts,amsmath,latexsym,amssymb} % these packages are required
% \usepackage{theorem}                % please use one of this two options for theorems
\usepackage{amsthm}                % please use one of this two options for theorems
\usepackage{mathrsfs,upref}         % not so essential, but part of journal style
\usepackage{mathptmx}		    % Journal is printed with poscript fonts: 
	                            % this package is essential for exact line and page breaks
		                    % comment if you do not have this package
\usepackage{jmi}	            % journal style, comment only if you find some bug.

\usepackage[T1]{fontenc}
\usepackage{hyperref}
\usepackage[nameinlink,noabbrev]{cleveref}

\usepackage{xcolor}

\usepackage{enumitem}
\newtheorem{thm}{Theorem}
\newtheorem{lem}[thm]{Lemma}
\newtheorem{prop}{Proposition}
\newtheorem{cor}[thm]{Corollary}
\newtheorem{example}[thm]{Example}
\newtheorem{rem}{Remark}
\numberwithin{equation}{section}
\Crefname{prop}{Proposition}{Propositions}


\begin{document}

\title{Weighted Hardy Inequalities for Nested Averages}

\author{Ludovick Bouthat and Pierre-Olivier Parisé}

\address{Ludovick Bouthat, 2325 Rue de l'Université, Québec, QC G1V 0A6\\
\email{ludovick.bouthat.1@ulaval.ca}}

\address{Pierre-Olivier Parisé, 3351 Bd des Forges, Trois-Rivières, QC G8Z 4M3,\\
\email{pierre-olivier.parise@uqtr.ca}}

\CorrespondingAuthor{Ludovick Bouthat}

\thanks{POP was supported by an NSERC Discovery Grant No. RGPIN-2026-04740.}

\date{DD.MM.YYYY}                               % Please, write the date of submission

\keywords{Hardy's inequality, weighted Hardy operators, nested averages, integral inequalities, discrete weighted inequalities}

\subjclass{Primary 26D15; Secondary 47B38, 47G10, 46E30}



\begin{abstract}
We study a family of Hardy-type inequalities for weighted averages over nested subsets of a measure space. Given a partition of a measure space and a weight function $m$, we consider operators of the form
\[
f \mapsto \frac{1}{M_n}\int_{X^{(n)}} m(x)f(x)\,\mathrm{d}\mu(x),
\]
with additional weights on the resulting sequence of averages. In particular, we generalize an inequality obtained by Vincent and Sohani in \cite{VincentSohani2025} and characterize the boundedness in terms of the finiteness of a single testing quantity $\beta$. We also provide two-sided estimates for the best constant $C_{\mathrm{opt}}$, namely
\[
\beta \leq C_{\mathrm{opt}} \leq p^{1/q} (p')^{1/p'}\beta \leq 2\beta.
\]
Thus the characterization is never off by more than a factor of 2. We also develop a second approach, inspired by Broadbent's proof of Hardy's inequality, which gives a local sufficient condition that often provides sharper constants and recovers several important cases, including the classical weighted Hardy inequality.
\end{abstract}

\maketitle



\section{Introduction}

In his work on integral equations, Hilbert proved that if $(a_n)_{n \geq 1}$ is a sequence of non-negative terms in $\ell^2$, then the double series
\[
\sum_{m=1}^{\infty}\sum_{n=1}^{\infty}
    \frac{a_m a_n}{m+n}
\]
is convergent \cite{Hilbert1906}. More precisely, Hilbert showed that this series is bounded by a constant multiple of $\sum_{n=1}^{\infty} a_n^2$. The sharp form of this result, with best constant $\pi$, was later obtained by Schur \cite{Schur1911}.

It was in the search for a simpler proof of Hilbert's inequality that Hardy was led to what is now known as Hardy's inequality \cite{hardy1919notes}. Hardy observed that Hilbert's inequality could be derived from the convergence of the sequence of arithmetic means: if $(a_n)_{n \geq 1} \in \ell^2$, then
\[
\sum_{n=1}^{\infty} \left( \frac{a_1+\cdots+a_n}{n} \right)^{\!2} \leq 4\sum_{n=1}^{\infty} |a_n|^2.
\]
Although this result first appeared as an auxiliary step, the mathematical community quickly recognized that it was of independent interest. This insight marked the beginning of a long and influential line of research centered on inequalities controlling averages of sequences or functions by their original norms. See for example the papers of T. A. A. Broadbent \cite{MR1574000}, E. B. Elliott \cite{MR1574962}, K. Grandjot \cite{MR1574375}, G. H. Hardy \cite{hardy1925notes}, T. Kaluza and G. Szeg\"{o} \cite{MR1574830}, and K. Knopp \cite{MR1574165}.

After Hardy communicated the above result to Marcel Riesz, the latter gave an alternative proof, published in a paper of the former \cite{MR1544414}, which yielded the following $p$-version: for $p>1$, $a_n \geq 0$, and $(a_n)_{n \geq 1} \in \ell^p$, one has
\[
\sum_{n=1}^{\infty} \left( \frac{a_1+\cdots+a_n}{n} \right)^p \leq C_p \sum_{n=1}^{\infty} a_n^p .
\]
The optimal value of the constant was later shown to be $C_p = \big( \frac{p}{p-1} \big)^p$, a result attributed to Landau, who supplied Hardy with a proof of the sharpness of this constant \cite{MR1575105}. In \cite{MR1544414}, Hardy also recorded the corresponding integral formulation, without proof, which bounds the $L^p$-norm of the averaging operator
\[
f \mapsto \frac{1}{x}\int_a^x f(t)\,\mathrm{d}t
\]
by the $L^p$-norm of $f$.

Recently, Bouthat and Mashreghi studied the boundedness as an operator on $\ell^2$ of the \emph{$L$-matrices} \cite{MRtmp}, which are defined according to a sequence $(a_n)_{n\geq 1}$ by setting
\[
L =[a_n]=
\begin{pmatrix}
a_0 & a_1 & a_2 & a_3 & \cdots \\
a_1 & a_1 & a_2 & a_3 & \cdots \\
a_2 & a_2 & a_2 & a_3 & \cdots \\
a_3 & a_3 & a_3 & a_3 & \cdots \\[-2pt]
\vdots & \vdots & \vdots & \vdots & \ddots \\
\end{pmatrix}.
\]
Since then, several authors studied these matrices (see \cite{MR4364583,MR4335803,MR4509958,MR4375658}). In this process, the inequality 
\begin{equation}\label{eq - initial}
    \sum_{n=1}^{\infty} \frac{1}{4^n} \left| \sum_{j=1}^{4^n} a_j \right|^2
\leq
C \sum_{n=1}^{\infty} |a_n|^2
\end{equation}
naturally appeared \cite{MR4364583}. This motivated Bouthat, Mashreghi and Morneau-Guérin to ask if this inequality was a particular instance of a more general family of inequalities. 

More specifically, the authors observed that the form of \eqref{eq - initial} is very reminiscent of Hardy's inequality in the case $p=2$, where the ordinary arithmetic means are replaced by more general weighted means taken over nested subsets of a sequence. In the above examples, the subsets of sequences are $N_n = \{1,2,\dots,4^n\}$. More generally, one may consider a partition of $\mathbb{N}$
\[
\mathbb{N} = N_1 \cup N_2 \cup \cdots
\]
and define
\[
\mathbf{N}_n := N_1\cup\cdots\cup N_n, \qquad n \geq 1.
\]
In that case, it was shown in \cite{Bouthat2023} that if $(a_n)_{n \geq 1}$ is a sequence of complex numbers, $(m_n)_{n \geq 1}$ and $(M_n)_{n \geq 1}$ are two sequences of weights, $1<p<\infty$, and\vspace{-2pt}
\begin{equation}\label{E:3}
\rho := \sup_{n \geq 1} \, \left(w_n \sum_{k=n}^{\infty}\frac{(w_1+\cdots+w_k)^{p-1}}{M_k^p}\right)^{\!\!1/p} \!< \infty,\vspace{-2pt}
\end{equation}
where
\begin{equation*}
w_n := \left(\sum_{j \in N_n} m_j^{p'} \right)^{\!\!1/p'}\!, \qquad n \geq 1\vspace{-2pt}
\end{equation*}
and $p'=\frac{p}{p-1}$, then\vspace{-4pt}
\[
\left(\sum_{n=1}^{\infty} \bigg|\frac{1}{M_n}\sum_{j \in \mathbf{N}_n} m_ja_j \bigg|^p\right)^{\!\!1/p}
\!\!\leq \rho  \left(\sum_{n=1}^{\infty} |a_n|^p\right)^{\!\!1/p}.
\]

Following this work, in 2025, Vincent and Sohani generalized this inequality even more to integral operators. In particular they proved the following result.

\begin{thm}[\cite{VincentSohani2025}, Theorem 2]\label{thm - VincentSohani}
Let $A_1,A_2,\ldots$ be a partition of $\mathbb{R}^m$ such that each $A_i$ is measurable. Let $B_n = \cup_{i=1}^n A_i$ and let $p,q \in (1,\infty)$ be such that $\frac{1}{p}+\frac{1}{p'}=1$. For $g>0$ defined on $\mathbb{R}^m$, define \vspace{-2pt}
\[ 
w_i = \left( \int_{A_i} g(x)^{p'} \,\mathrm{d}x \right)^{\!\smash{1/p'}}. 
\] 
If $M_n$ is a sequence of positive numbers such that 
\[ 
\rho := \sup_{n \geq 1} \, \left(w_n \sum_{k=n}^{\infty}\frac{(w_1+\cdots+w_k)^{p-1}}{M_k^p}\right)^{\!\!1/p} < \infty, 
\] 
then 
\[ 
\left( \sum_{n=1}^{\infty} \left| \frac{1}{M_n} \int_{B_n} g(x) f(x) \,\mathrm{d}x \right|^p \right)^{\!\!1/p} \!\leq \rho \left( \int_{\mathbb{R}^m} |f(x)|^p \,\mathrm{d}x \right)^{\!\!1/p}. 
\] 
\end{thm}


While Vincent and Sohani considerably generalized the inequality of Bouthat, Mashreghi and Morneau-Guérin, there are still some limitations to both of them. For instance, one can show that, while both of these inequalities look similar to the classical Hardy inequality, neither of them allows to recover this simple case. Moreover, while they are enough to obtain \eqref{eq - initial} with a certain constant $C>0$, they are not sufficiently precise in this regime to identify the sharp multiplicative constant.

More generally, neither of the sufficient conditions for these inequalities is also necessary. Lastly, one may notice that there is still room for apparent generalization. In particular, one may consider a general measure space $(X,\mu)$, consider mixed norms, or add another sequence of weights in the inequality, which would allow to recover more past generalizations of Hardy's inequality in the literature (see \Cref{sec - local}). Note that here and throughout, a sequence of weights is meant to be any sequence of positive numbers. It is the objective of this paper to simultaneously address all of these concerns.

Hence, the main question considered in this paper is the following: \emph{Let $1 \leq p\leq q\leq \infty$. Let $(X,\mu)$ be a measure space. Let $X = X_1 \cup X_2 \cup \cdots$ be a partition of $X$ and set $X^{(n)} := X_1 \cup \cdots \cup X_n$. Let $(b_n)_{n\geq 1}$ and $(M_n)_{n\geq 1}$ be sequences of weights. Let $m$ be a positive measurable function on $X$ and $f$ a complex measurable function on $X$. Under which conditions is there a number $\rho>0$ such that
\begin{equation}\label{eq - goal}
    \left(\sum_{n = 1}^\infty b_n \left| \frac{1}{M_n} \int_{X^{(n)}} m (x) f (x)  \mathrm{d}\mu (x) \right|^q\right)^{\!\!1/q} \!\!\leq \rho \left(\sum_{n = 1}^\infty b_n \int_{X_n} |f (x)|^p  \mathrm{d}\mu (x) \right)^{\!\!1/p}.
\end{equation}
}

The paper is outlined as follows. In \Cref{sec - charac}, we begin by showing that, while the generalizations proposed above are convenient and interesting, they can be equivalently reduced to the discrete case. As a corollary, we are able to completely characterize the condition under which a constant $\rho>0$ exists such that \eqref{eq - goal} holds true by using a 1978 result of Bradley \cite{Bradley}. Some examples are then studied to compare the sharpness of the constant $\rho$ obtained from this approach to the constants obtained in \cite{Bouthat2023} and \cite{VincentSohani2025}.

While this proposed approach provides a complete characterization, we shall see that the constants obtained are often not sharp. Consequently, in \Cref{sec - local}, we will explore another approach which gives a local sufficient condition for the validity of \eqref{eq - goal}. While the locality of the condition is a strong hypothesis to satisfy, we shall see that it often comes with the benefit of giving some sharper constants.

Lastly, we conclude in \Cref{sec - conclusion} by raising some open questions and research avenues for these kinds of inequalities.



\section{A complete characterization}\label{sec - charac}

As outlined above, we first prove an equivalence between the measure-theoretic case and the discrete case. Note that here and in the rest of the paper, we write $p'$ to denote the Hölder conjugate of $p$, i.e., the number satisfying $\frac{1}{p}+\frac{1}{p'}=1$, and we define
\begin{equation}\label{def - wn}
    w_n := \begin{cases}
        \left( \int_{X_n} |m(x)|^{p'}\,\mathrm{d}\mu(x) \right)^{\!\smash{1/p'}} \quad&\text{if }p>1,\\
        %
        \operatorname*{ess\,sup}_{x \in X_n} |m(x)|\quad&\text{if }p=1.
    \end{cases}
\end{equation}

% we begin with the following results, which allows us to reduce the problem to the simpler discrete case. We split the result into two cases : when $1 < p < \infty$ and $p = 1$.

% \subsection{Case $1 < p < \infty$}
% As promised, we show that the general problem reduces to the simpler discrete case.

\begin{prop}\label{prop - equiv}
Let $1\leq p\leq q\leq\infty$. Let $(X,\mu)$ be a measure space, with $\mu$ semi-finite if $p=1$. Let $X = X_1 \cup X_2 \cup \cdots$ be a partition of $X$ and set $X^{(n)} := X_1 \cup \cdots \cup X_n$. Let $(b_n)_{n\geq 1}$ and $(M_n)_{n\geq 1}$ be sequences of weights. Let $m$ be an $L^{p'} (X)$ function and define $w_n$ as in \eqref{def - wn}. Then there exists $\rho>0$ such that the inequality
\begin{equation}\label{eq - measure-hardy}
    \left(\sum_{n=1}^{\infty} b_n \left| \frac{1}{M_n} \int_{X^{(n)}}m(x)f(x)\,\mathrm{d}\mu(x) \right|^q\right)^{\!\!1/q}\! \leq \rho \left(\sum_{n=1}^{\infty} b_n \int_{X_n}|f(x)|^p\,\mathrm{d}\mu(x) \right)^{\!\!1/p}
\end{equation}
holds for every measurable function $f$ if and only if the discrete inequality
\begin{equation}\label{eq - discrete-hardy}
    \left(\sum_{n=1}^{\infty} b_n \bigg|\frac{1}{M_n} \sum_{k=1}^{n}w_k a_k \bigg|^q \right)^{\!\!1/q}\! \leq \rho \left(\sum_{k=1}^{\infty} b_k |a_k|^p\right)^{\!\!1/p}
\end{equation}
holds for every sequence $a=(a_k)_{k\geq 1}$.% \!$($\!with the obvious modifications if $p,q=\infty)$.
\end{prop}


\begin{rem}\label{rem - explain}
    Here, the case $p=\infty$ and $q=\infty$ should be interpreted in the natural way. Namely, the left and right-hand sides of \eqref{eq - measure-hardy} should be understood respectively as %\textbf{[Je ne suis pas sûr de comprendre pourquoi le $b_n$ disparait.]} \textcolor{red}{[Parce que si on l'inclu dans la valeur absolue $|\cdot|^q$ (ou $|\cdot|^p$), on doit ajouter une puissance $b_n^{1/q}$ ou $b_n^{1/p}$. Quand $q$ ou $p$ tend vers l'infini, on obtient donc essentiellement le supremum du truc dans la valeur absolue, mais en même temps on a $b_n^{1/q}\to 1$ si $b_n\neq 0$. Donc le $b_n$ disparait et on garde juste le suprémum sur les indices où $b_n\neq 0$.]}
    \[
    \sup_{b_n>0} \left| \frac{1}{M_n} \int_{X^{(n)}}m(x)f(x)\,\mathrm{d}\mu(x) \right| \qquad\text{and}\qquad \operatorname*{ess\,sup}_{\cup_{b_n> 0} X_n}|f(x)|,
    \]
    while the left and right-hand sides of \eqref{eq - discrete-hardy} should be interpreted respectively as
    \[
    \sup_{b_n>0} \bigg|\frac{1}{M_n} \sum_{k=1}^{n}w_k a_k \bigg| \qquad\text{and}\qquad \sup_{b_n>0}|a_n|.
    \]
\end{rem}

\begin{proof}
We first prove that \eqref{eq - discrete-hardy} implies \eqref{eq - measure-hardy}. For each $k\geq 1$, define
% \[
% F_k:= \left( \int_{X_k}|f(x)|^p\,\mathrm{d}\mu(x) \right)^{\!\!1/p}.
% \]
\[
F_k:= 
\begin{cases}
    \left( \int_{X_k} |f(x)|^{p}\,\mathrm{d}\mu(x) \right)^{\!\smash{1/p}} \quad&\text{if }p<\infty,\\
    %
    \operatorname*{ess\,sup}_{x \in X_k} |f(x)|\quad&\text{if }p=\infty.
\end{cases}
\]
By Hölder's inequality,
\[
\left| \int_{X_k}m(x)f(x)\,\mathrm{d}\mu(x) \right| \leq \left( \int_{X_k}|m(x)|^{p'}\,\mathrm{d}\mu(x) \right)^{\!\!1/p'} \!\left( \int_{X_k}|f(x)|^p\,\mathrm{d}\mu(x) \right)^{\!\!1/p} \!= w_k F_k,
\]
with the obvious modifications in the inequality if $p=1$ or $\infty$. Therefore,
\[
\left| \int_{X^{(n)}}m(x)f(x)\,\mathrm{d}\mu(x) \right| = \left| \sum_{k=1}^{n} \int_{X_k}m(x)f(x)\,\mathrm{d}\mu(x) \right| \leq \sum_{k=1}^{n}w_k F_k .
\]
Applying \eqref{eq - discrete-hardy} to the sequence $(F_k)_{k\geq 1}$ gives
\begin{align*}
    \left(\sum_{n=1}^{\infty} b_n \bigg| \frac{1}{M_n} \int_{X^{(n)}}m(x)f(x)\,\mathrm{d}\mu(x) \bigg|^q\right)^{\!\!1/q}\!\!
    &\leq \left(\sum_{n=1}^{\infty} b_n \bigg| \frac{1}{M_n} \sum_{k=1}^{n}w_k F_k \bigg|^q\right)^{\!\!1/q}\\
    %
    &\leq \rho \!\left( \sum_{k=1}^{\infty}b_k F_k^p \right)^{\!\!1/p}\!, %\\
    %
    % &= \rho \left(\sum_{k=1}^{\infty}b_k \int_{X_k}|f(x)|^p\,\mathrm{d}\mu(x) \right)^{\!\!1/p},
\end{align*}
which is \eqref{eq - measure-hardy}.

Conversely, suppose that \eqref{eq - measure-hardy} holds. We prove \eqref{eq - discrete-hardy}. If $w_k=0$, then the $k$th coordinate does not contribute to the left-hand side of \eqref{eq - discrete-hardy}; hence we may assume $a_k=0$ whenever $w_k=0$. We write $W := \{ n \in \mathbb{N} \, : \, w_n > 0 \}$. 


Let $k \in W$. If $p=1$, then $w_k=\operatorname*{ess\,sup}_{X_k}|m|$. Hence, since $\mu$ is semi-finite, choose a set $E_{k,\varepsilon}\subset X_k$ of finite positive measure on which $|m|>w_k-\varepsilon$. Now define $f_{k}$ on $X_k$ by
\[
f_{k} := \begin{cases} 
    a_k e^{-i \arg(m)} \frac{\chi_{E_{k,\varepsilon}} }{\mu(E_{k,\varepsilon})} & \text{ if } p = 1, \\
    a_k w_k^{1-p'}e^{-i \arg(m)}|m|^{p'-1} \chi_{X_k} & \text{ if } p > 1,    
    \end{cases}
\qquad 
\]
where $\arg(m)$ is arbitrary when $m=0$, and set $f := \sum_{k \in W} f_{k}$. If $p=1$, then
\[
\int_{X_k}|f(x)|^p\,\mathrm{d}\mu(x) = \int_{X_k}|a_k|\frac{\chi_{E_{k,\varepsilon}}}{\mu(E_{k,\varepsilon})}\,\mathrm{d}\mu(x) = |a_k|.
\]
If $1<p<\infty$, then
\[
\begin{aligned}
    \int_{X_k}|f(x)|^p\,\mathrm{d}\mu(x)
    %&= |a_k|^p w_k^{p(1-p')} \int_{X_k}|m(x)|^{p(p'-1)}\,\mathrm{d}\mu(x) \\
    %
    &= |a_k|^p w_k^{p(1-p')} \int_{X_k}|m(x)|^{p'}\,\mathrm{d}\mu(x) %\\
    %
    % &= |a_k|^p w_k^{p(1-p')}w_k^{p'} ~=~ 
    = |a_k|^p .
\end{aligned}
\]
Hence, in both cases, we have
\[
\left(\sum_{n=1}^{\infty} b_n \int_{X_n}|f(x)|^p\,\mathrm{d}\mu(x) \right)^{\!\!1/p} = \left(
\sum_{k=1}^{\infty} b_k |a_k|^p \right)^{\!\!1/p}.
\]
Lastly, if $p=\infty$, then we verify directly that
\[
\left(\sum_{n=1}^{\infty} b_n \int_{X_n}|f(x)|^p\,\mathrm{d}\mu(x) \right)^{\!\!1/p} =\operatorname*{ess\,sup}_{\cup_{b_n> 0} X_n}|f(x)| = \sup_{b_n> 0}|a_n|
= \left(\sum_{k=1}^{\infty} b_k |a_k|^p \right)^{\!\!1/p}.
\]
Moreover, we always have
\[
\int_{X_k}m(x)f(x)\,\mathrm{d}\mu(x) = \lambda_{k,\varepsilon}a_k,
\]
where $w_k-\varepsilon \leq \lambda_{k,\varepsilon} \leq w_k$ if $p=1$ and $\lambda_{k,\varepsilon}=w_k$ if $p>1$. In the case $p = 1$, applying \eqref{eq - measure-hardy}, we find 
\begin{align*}
    \left( \sum_{n=1}^{\infty} b_n \bigg| \frac1{M_n} \sum_{k=1}^{n} \lambda_{k,\varepsilon}a_k \bigg|^q \right)^{\!\!1/q}\!\! &= \left(\sum_{n=1}^{\infty} b_n \left| \frac{1}{M_n} \int_{X^{(n)}}m(x)f(x)\,\mathrm{d}\mu(x) \right|^q\right)^{\!\!1/q} \\
    %
    &\leq \rho \!\left(\sum_{n=1}^{\infty} b_n \int_{X_n}|f(x)|^p\,\mathrm{d}\mu(x) \right)^{\!\!1/p} \!\!= \rho\!\left(\sum_{k=1}^{\infty} b_k |a_k|^p \right)^{\!\!1/p}\!.
\end{align*}
The result then follows directly by letting $\varepsilon\to0$ and using Fatou's lemma. Note that the case $q=\infty$ is handled in the same way, with the weighted $\ell^q$-norm replaced by the supremum norm on the indices satisfying $b_n>0$.
%
% Then
% \[
% \int_{X_k}m(x)f(x)\,\mathrm{d}\mu(x) = a_kw_k^{1-p'}
% \int_{X_k}|m(x)|^{p'}\,\mathrm{d}\mu(x) = a_kw_k .
% \]
% Moreover, since $p(p'-1)=p'$, we have
% \[
% \begin{aligned}
%     \int_{X_k}|f(x)|^p\,\mathrm{d}\mu(x)
%     &= |a_k|^p w_k^{p(1-p')} \int_{X_k}|m(x)|^{p(p'-1)}\,\mathrm{d}\mu(x) \\
%     %
%     &= |a_k|^p w_k^{p(1-p')} \int_{X_k}|m(x)|^{p'}\,\mathrm{d}\mu(x) \\
%     %
%     &= |a_k|^p w_k^{p(1-p')}w_k^{p'} ~=~ |a_k|^p .
% \end{aligned}
% \]
% Applying \eqref{eq - measure-hardy} to this function $f$, we obtain
% \[
% \left(\sum_{n=1}^{\infty} b_n
% \bigg|
%     \frac{1}{M_n}
%     \sum_{k=1}^{n}w_k a_k
% \bigg|^q\right)^{\!\!1/q}\!
% \leq
% \rho \left(
% \sum_{k=1}^{\infty} b_k |a_k|^p \right)^{\!\!1/p}.
% \]
% This is precisely \eqref{eq - discrete-hardy}.
%
% Now consider the case $p=1$. By the triangle inequality and a standard truncation argument, it is enough to prove \eqref{eq - discrete-hardy} for finitely supported sequences $a=(a_k)_{k\ge1}$. Let $a$ be finitely supported. We may assume that $w_k>0$ whenever $a_k\neq0$. Fix $\varepsilon>0$. Since
% \[
% w_k=\operatorname*{ess\,sup}_{x\in X_k}|m(x)|,
% \]
% the set
% \[
% \{x\in X_k:m(x)>w_k-\varepsilon\}
% \]
% has positive measure. Choose a measurable set $E_k\subset X_k$ such that
% \[
% 0<\mu(E_k)<\infty
% \qquad\text{and}\qquad
% m(x)>w_k-\varepsilon
% \quad\text{for }x\in E_k,
% \]
% and define
% \[
% f(x):=\sum_k \frac{a_k}{\mu(E_k)}\chi_{E_k}(x),
% \]
% where the sum is over the finite support of $a$. Then
% \[
% \int_{X_k}|f(x)|\,\mathrm{d}\mu(x)=|a_k|
% \]
% and
% \[
%     \int_{X_k}m(x)f(x)\,d\mu(x)
%     =
%     \lambda_k a_k,
%     \qquad
%     \lambda_k:=\frac{1}{\mu(E_k)}\int_{E_k}m(x)\,d\mu(x).
% \]
% Moreover,
% \[
%     w_k-\varepsilon\leq \lambda_k\leq w_k .
% \]
% Applying (2.2) to this function $f$, we get
% \[
%     \left(
%     \sum_{n=1}^{\infty}
%     b_n
%     \left|
%     \frac1{M_n}\sum_{k=1}^n \lambda_k a_k
%     \right|^q
%     \right)^{1/q}
%     \leq
%     \rho
%     \sum_{k=1}^{\infty}b_k|a_k|.
% \]
% Letting $\varepsilon\to0$, we obtain, for every finitely supported sequence
% $a$,
% \[
%     \left(
%     \sum_{n=1}^{\infty}
%     b_n
%     \left|
%     \frac1{M_n}\sum_{k=1}^n w_k a_k
%     \right|^q
%     \right)^{1/q}
%     \leq
%     \rho
%     \sum_{k=1}^{\infty}b_k|a_k|.
% \]
% The result for arbitrary sequences follows by applying this estimate to the truncations $a^{(N)}=(a_1,\dots,a_N,0,0,\dots)$ and letting $N\to\infty$, using Fatou's lemma. The case $q=\infty$ is obtained in the same way, replacing the $L^q$-norm by the supremum norm.
\end{proof}


The preceding proposition shows that the measure-theoretic problem considered above is equivalent to a discrete  inequality. The corresponding continuous problem has a classical answer, due to Muckenhoupt \cite{Muckenhoupt}. More precisely, if $1<p<\infty$, then
\[
\left(\int_0^\infty \left|u(x)\int_0^x f(t)\,\mathrm{d}t \right|^p \,\mathrm{d}x \right)^{\!\!1/p}\!
\le C \left( \int_0^\infty |v(x)f(x)|^p\,\mathrm{d}x \right)^{\!\!1/p}
\]
holds for all measurable functions $f$ if and only if
\[
A_p:= \sup_{s>0} \left(\int_s^\infty |u(x)|^p\,\mathrm{d}x\right)^{\!\!1/p}\! \left(\int_0^s |v(x)|^{-p'}\,\mathrm{d}x\right)^{\!\!1/p'} <\infty .
\]
Moreover, if $C_{\mathrm{opt}}$ is the best possible constant, then $A_p\le C_{\mathrm{opt}} \le p^{1/p}(p')^{1/p'}A_p .$ 
%
Bradley \cite{Bradley} later extended this result to mixed norms. Namely, for $1\le p\le q\le\infty$, the inequality\vspace{-1pt}
\[
\left(\int_0^\infty\left|u(x)\int_0^x f(t)\,\mathrm{d}t\right|^q \mathrm{d}x\right)^{\!\!1/q} \!\le C\left(\int_0^\infty |v(x)f(x)|^p\,\mathrm{d}x\right)^{\!\!1/p}
\]
is characterized by the same type of testing condition. That is,\vspace{-1pt}
\[
A_{p,q}:= \sup_{s>0} \left(\int_s^\infty |u(x)|^q\,\mathrm{d}x\right)^{\!\!1/q} \!\left(\int_0^s |v(x)|^{-p'}\,\mathrm{d}x\right)^{\!\!1/p'}\! <\infty ,
\]
with the usual modifications when $p=1$ or $q=\infty$. In particular, the case $q=p$ recovers Muckenhoupt's theorem.

The results of Muckenhoupt and Bradley are formulated for integral operators on $(0,\infty)$. In the present paper, however, \Cref{prop - equiv} reduces the problem to a discrete operator. The corresponding discrete criterion may be obtained from the continuous one by a standard step-function discretization, or by repeating the proof with sums in place of integrals. We state this discrete form explicitly in the notation needed below.

\begin{prop}[Discrete Muckenhoupt--Bradley criterion]\label{prop - old}
Let $1\leq p\leq q\leq \infty$, and let
$(u_n)_{n\ge1}$ and $(v_n)_{n\ge1}$ be sequences of weights. Then there exists some $C>0$ such that\vspace{-1pt}
\[
\left(\sum_{n=1}^{\infty} \bigg| u_n\sum_{k=1}^n a_k \bigg|^q \right)^{\!\!1/q}\!\! \le C \left(\sum_{k=1}^{\infty}|v_ka_k|^p \right)^{\!\!1/p}
\]
holds for every sequence $a=(a_k)_{k\ge1}$ if and only if\vspace{-1.5pt}
\[
\beta := \sup_{N\ge1} \left(\sum_{n=N}^{\infty} |u_n|^q\right)^{\!\!1/q}\! \left(\sum_{k=1}^N v_k^{-p'}\right)^{\!\!1/p'}\!
<\infty
\]
\!$($\!with the obvious modifications if $p,q\in\{1,\infty\}$, as in \Cref{rem - explain}\!$)$. Moreover, if $C_{\mathrm{opt}}$ denotes the best possible constant, then $\beta \le C_{\mathrm{opt}} \le p^{1/q}(p')^{1/p'}\beta .$
\end{prop}

\begin{rem}
    If $p=1$, then $p^{1/q}(p')^{1/p'}=1$. Hence, the optimal constant is \vspace{-1.5pt}
    \begin{equation}\label{eq - beta}
        \beta = \sup_{N\ge1} \max_{1\leq k \leq N} \frac{1}{v_k}\left(\sum_{n=N}^{\infty} |u_n|^q\right)^{\!\!1/q}.
    \end{equation}
    However, this constant may be replaced with an equivalent quantity, namely\vspace{-1.5pt}
    $$
        \beta' := \sup_{N \geq 1} \frac{1}{v_N} \left( \sum_{n = N}^\infty |u_n|^q \right)^{1/q}, 
    $$
    which is easier to handle in applications than \eqref{eq - beta}. Indeed, the inequality $\beta'\le\beta$ is immediate. Conversely, if $k\le N$, then $\left(\sum_{n=N}^{\infty} |u_n|^q\right)^{1/q}\le \left(\sum_{n=k}^{\infty} |u_n|^q\right)^{1/q}$, and therefore
    \[
    \frac{1}{v_k}\left(\sum_{n=N}^{\infty} |u_n|^q\right)^{\!\!1/q}
    \le
    \frac{1}{v_k}\left(\sum_{n=k}^{\infty} |u_n|^q\right)^{\!\!1/q}
    \le
    \beta' .
    \]
    Maximizing over $1\le k\le N$ and then over $N$ yields $\beta\le\beta'$. Hence $\beta=\beta'$.
\end{rem}


% In particular, applying this proposition with
% \[
% u_n=\frac{b_n^{1/q}}{M_n}
% \qquad\text{and}\qquad
% v_n=b_n^{1/p},
% \]
% and replacing $a_n$ by $w_na_n$ gives the characterization required for the inequality obtained in \Cref{prop - equiv}.


% The result above shows that the general case is equivalent to the discrete one. Therefore, we focus on the discrete case, for which a characterization of boundedness is provided for the case $1 < p < \infty$. Note that the following result is a standard two-weight discrete Hardy inequality, going back in continuous form to Muckenhoupt \cite{Muckenhoupt} and in discrete form to Andersen and Heinig \cite{AndersenHeinig}. We include the proof, in the present notation, to make clear how the weights arising from the nested-average problem interact with the testing quantity. \textcolor{red}{[Should we include the full proof or just how we can get it by a change of variable of the known inequality? Démontrer sans le $w_k$.]}

% {\color{purple}
% \begin{thm}\label{thm - charac}
% Let $1\leq p\leq q\leq \infty$ and $p'=p/(p-1)$. Let $(u_n)_{n\geq 1}$, $(v_n)_{n\geq 1}$, and $(w_n)_{n\geq 1}$ be positive sequences. Set
% \[
% \beta := \sup_{N\geq 1} \left(\sum_{n=N}^{\infty}u_n\right)^{\!\!1/q} \!\left(\sum_{k=1}^{N} w_k^{p'}v_k^{-p'/p} \right)^{\!\!1/p'}.
% \]
% Then the inequality
% \[
% \left(\sum_{n=1}^{\infty}u_n \left|\sum_{k=1}^{n}w_k a_k\right|^q \right)^{\!\!1/q} \leq C \left(\sum_{k=1}^{\infty}v_k |a_k|^p\right)^{\!\!1/p}
% \]
% holds for all sequences $a=(a_k)_{k\geq 1}$ and some $C>0$ if and only if $\beta<\infty$. Moreover, if $C_{\mathrm{opt}}$ denotes the best possible constant, then
% \[
% \beta \leq C_{\mathrm{opt}} \leq p^{1/q}(p')^{1/p'}\beta .
% \]
% \end{thm}

% \begin{proof}
% It is enough to prove the result for nonnegative sequences. The general case follows simply by replacing $a_k$ with $|a_k|$. For $n\geq 1$, define
% \[
% W_n:=w_n^{p'} v_n^{-p'/p},
% \qquad
% V_n:=\sum_{k=1}^{n}W_k
% \quad\text{ and }\quad
% R_n:=\sum_{j=n}^{\infty}u_j .
% \]
% Then the condition $\beta<\infty$ is equivalent to $R_n^{1/q}V_n^{1/p'}\leq \beta$ for every $n\geq 1$, or equivalently
% \begin{equation}\label{eq - bound}
%     R_n\leq \beta^q V_n^{-q/p'}.
% \end{equation}

% We first prove necessity. Suppose that the inequality holds with constant $C$.
% Fix $N\geq 1$ and define
% \[
% a_k=
% \begin{cases}
%     w_k^{p'-1}v_k^{-p'/p}, & k\leq N,\\
%     0, & k>N.
% \end{cases}
% \]
% Then, for every $n\geq N$, we have
% \[
% \sum_{k=1}^{n}w_k a_k=V_N.
% \]
% Therefore,
% \[
% \sum_{n=1}^{\infty}u_n \!\left(\sum_{k=1}^{n}w_k a_k\right)^{\!q} = \sum_{n=1}^{N-1}u_n \!\left(\sum_{k=1}^{n}w_k a_k\right)^{\!q} +  V_N^q\sum_{n=N}^{\infty}u_n \geq V_N^q\sum_{n=N}^{\infty}u_n .
% \]
% Moreover, since $p(p'-1)=p'$, we have
% \[
% \sum_{k=1}^{\infty}v_k a_k^p = \sum_{k=1}^{N} v_k w_k^{p(p'-1)}v_k^{-p'} = \sum_{k=1}^{N}w_k^{p'} v_k^{-p'/p} = V_N .
% \]
% Using the assumed inequality, we obtain
% \[
% \left(V_N^q\sum_{n=N}^{\infty}u_n\right)^{\!\!1/q} \leq C V_N^{1/p}.
% \]
% Thus
% \[
% \left(\sum_{n=N}^{\infty}u_n\right)^{\!\!1/q} \!V_N^{1/p'} \leq C < \infty.
% \]
% Taking the supremum over $N$ gives $\beta\leq C$. Hence
% \[
% \beta\leq C_{\mathrm{opt}} .
% \]

% We now prove sufficiency. Assume that $\beta<\infty$. Let
% $0<\theta<1$. For each $n\geq 1$, we write
% \[
% \sum_{k=1}^{n}w_k a_k = \sum_{k=1}^{n} \left(v_k^{1/p}a_k V_k^{\theta/p'}\right) \!\left(w_k v_k^{-1/p}V_k^{-\theta/p'}\right).
% \]
% By Hölder's inequality,
% \[
% \left(\sum_{k=1}^{n}w_k a_k\right)^{\!q} \leq \left( \sum_{k=1}^{n}v_k a_k^p V_k^{\theta(p-1)} \right)^{\!q/p}\! \left( \sum_{k=1}^{n}w_k^{p'} v_k^{-p'/p}V_k^{-\theta} \right)^{\!q/p'}.
% \]
% By definition, $W_k=w_k^{p'} v_k^{-p'/p}$ and $V_k=\sum_{j=1}^{k}W_j$. Hence, $W_k = V_k-V_{k-1}$ and
% \begin{align*}
%     \sum_{k=1}^{n}W_k V_k^{-\theta} &= \sum_{k=1}^{n}(V_k-V_{k-1}) V_k^{-\theta} = \sum_{k=1}^{n}\int_{V_{k-1}}^{V_k}V_k^{-\theta}\,dt \\
%     %
%     &\leq  \sum_{k=1}^{n}\int_{V_{k-1}}^{V_k} t^{-\theta}\,dt = \int_0^{V_n} t^{-\theta}\,dt = \frac{V_n^{1-\theta}}{1-\theta}.
% \end{align*}
% Therefore,
% \[
% \left(\sum_{k=1}^{n}w_k a_k\right)^{\!p} \leq \frac{V_n^{(1-\theta)(p-1)}}{(1-\theta)^{p-1}} \sum_{k=1}^{n}v_k a_k^p V_k^{\theta(p-1)} .
% \]
% Multiplying by $u_n$ and summing over $n$, we get
% \begin{align*}
%     \sum_{n=1}^{\infty}u_n \!\left(\sum_{k=1}^{n}w_k a_k\right)^{\!p}
%     &\leq \frac{1}{(1-\theta)^{p-1}} \sum_{n=1}^{\infty} u_n V_n^{(1-\theta)(p-1)} \sum_{k=1}^{n}v_k a_k^p V_k^{\theta(p-1)} \\
%     %
%     &= \frac{1}{(1-\theta)^{p-1}} \sum_{k=1}^{\infty} v_k a_k^p V_k^{\theta(p-1)} \sum_{n=k}^{\infty} u_n V_n^{(1-\theta)(p-1)} .
% \end{align*}
% Let $\gamma:=(1-\theta)(p-1)$. We claim that
% \begin{equation}\label{eq - bound2}
%     \sum_{n=k}^{\infty}u_nV_n^{\gamma} \leq \frac{\beta^p}{\theta} V_k^{-\theta(p-1)} .
% \end{equation}
% Indeed, by summation by parts and by \eqref{eq - bound},
% \[
% \begin{aligned}
%     \sum_{n=k}^{\infty}u_nV_n^\gamma &= \sum_{n=k}^{\infty} (R_{n}-R_{n+1})V_n^\gamma = \sum_{n=k}^{\infty} R_{n}V_n^\gamma - \sum_{n=k}^{\infty} R_{n+1}V_n^\gamma \\
%     %
%     &= R_{k}V_k^\gamma + \sum_{n=k+1}^{\infty} R_{n}V_n^\gamma - \sum_{n=k+1}^{\infty} R_{n}V_{n-1}^\gamma \\
%     %
%     &\leq \beta^p V_k^{\gamma+1-p} + \beta^p \sum_{n=k+1}^{\infty} V_n^{1-p} \!\left(V_n^\gamma-V_{n-1}^\gamma\right).
% \end{aligned}
% \]
% Since $1-p<0$, the last sum is bounded above by the corresponding integral:
% \begin{align*}
%     \sum_{n=k+1}^{\infty} V_n^{1-p} \!\left(V_n^\gamma-V_{n-1}^\gamma\right) &= \sum_{n=k+1}^{\infty} \gamma\int_{V_{n-1}}^{V_n} V_n^{1-p}t^{\gamma-1} \,dt \leq \sum_{n=k+1}^{\infty} \gamma\int_{V_{n-1}}^{V_n} t^{1-p}t^{\gamma-1} \,dt \\
%     %
%     &= \gamma\int_{V_{k}}^{V_\infty} t^{\gamma-p} \,dt =\frac{\gamma}{p-1-\gamma}\bigl(V_k^{\gamma+1-p} - V_\infty^{\gamma+1-p} \bigr) \\
%     %
%     &\leq  \frac{\gamma}{p-1-\gamma} V_k^{\gamma+1-p},
% \end{align*}
% where we used the fact that $V_\infty:=\lim_{n\to\infty} V_n \geq 0$ and $p-1-\gamma=\theta(p-1)>0$.
% % \[
% %     \sum_{n=k+1}^{\infty}
% %     V_n^{1-p}
% %     \left(V_n^\gamma-V_{n-1}^\gamma\right)
% %     \leq
% %     \int_{V_k}^{\infty} t^{1-p}\,d(t^\gamma).
% % \]
% % Because $\gamma<p-1$, we have
% % \[
% %     \int_{V_k}^{\infty} t^{1-p}\,d(t^\gamma)
% %     =
% %     \gamma\int_{V_k}^{\infty}t^{\gamma-p}\,dt
% %     =
% %     \frac{\gamma}{p-1-\gamma}V_k^{\gamma+1-p}.
% % \]
% Hence
% \[
% \begin{aligned}
%     \sum_{n=k}^{\infty}u_nV_n^{\gamma}
%     &\leq
%     \beta^p\left(1+\frac{\gamma}{p-1-\gamma}\right) V_k^{\gamma+1-p} = \frac{p-1}{p-1-\gamma}\,\beta^p\,V_k^{\gamma+1-p}.
% \end{aligned}
% \]
% Since $\gamma=(1-\theta)(p-1)$, we have
% \[
% \frac{p-1}{p-1-\gamma} = \frac{1}{\theta} \quad\text{ and }\quad \gamma+1-p=-\theta(p-1),
% \]
% proving \eqref{eq - bound2}. Substituting \eqref{eq - bound2} into the previous estimate gives
% \[
% \begin{aligned}
%     \sum_{n=1}^{\infty}u_n\!\left(\sum_{k=1}^{n}w_k a_k\right)^{\!p} &\leq \frac{\beta^p}{\theta(1-\theta)^{p-1}} \sum_{k=1}^{\infty} v_k a_k^p V_k^{\theta(p-1)} V_k^{-\theta(p-1)} \\
%     %
%     &= \frac{\beta^p}{\theta(1-\theta)^{p-1}} \sum_{k=1}^{\infty}v_k a_k^p .
% \end{aligned}
% \]
% Optimizing over $0<\theta<1$, we find that the minimum is attained at $\theta=1/p$, where
% \[
% \frac{1}{\theta(1-\theta)^{p-1}} = p q^{p-1}.
% \]
% Consequently,
% \[
% \sum_{n=1}^{\infty}u_n \!\left(\sum_{k=1}^{n}w_k a_k\right)^{\!p} \leq p q^{p-1}\beta^p \sum_{k=1}^{\infty}v_k a_k^p .
% \]
% Taking $p$-th roots yields
% \[
% C_{\mathrm{opt}} \leq p^{1/p}q^{1/q}\beta .
% \]
% Together with the necessity proved above, this completes the proof.
% \end{proof}
% }

% \begin{thm}\label{thm - charac}
% Let $1<p<\infty$ and let $q=p/(p-1)$. Let $(u_n)_{n\geq 1}$, $(v_n)_{n\geq 1}$, and $(w_n)_{n\geq 1}$ be positive sequences. Set
% \[
% \beta := \sup_{N\geq 1} \left(\sum_{n=N}^{\infty}u_n\right)^{\!\!1/p} \!\left(\sum_{k=1}^{N}w_k^q v_k^{-q/p}\right)^{\!\!1/q}.
% \]
% Then the inequality
% \[
% \sum_{n=1}^{\infty}u_n \left|\sum_{k=1}^{n}w_k a_k\right|^p \leq C^p \sum_{k=1}^{\infty}v_k |a_k|^p
% \]
% holds for all sequences $a=(a_k)_{k\geq 1}$ and some $C>0$ if and only if
% $\beta<\infty$. Moreover, if $C_{\mathrm{opt}}$ denotes the best possible
% constant, then
% \[
% \beta \leq C_{\mathrm{opt}} \leq p^{1/p}q^{1/q}\beta .
% \]
% \end{thm}

% \begin{proof}
% It is enough to prove the result for nonnegative sequences. The general case follows simply by replacing $a_k$ with $|a_k|$. For $n\geq 1$, define
% \[
% W_n:=w_n^q v_n^{-q/p},
% \qquad
% V_n:=\sum_{k=1}^{n}W_k
% \quad\text{ and }\quad
% R_n:=\sum_{j=n}^{\infty}u_j .
% \]
% Then the condition $\beta<\infty$ is equivalent to $R_n^{1/p}V_n^{1/q}\leq \beta$ for every $n\geq 1$, or equivalently
% \begin{equation}\label{eq - bound}
%     R_n\leq \beta^p V_n^{1-p}.
% \end{equation}

% We first prove necessity. Suppose that the inequality holds with constant $C$.
% Fix $N\geq 1$ and define
% \[
%     a_k=
%     \begin{cases}
%         w_k^{q-1}v_k^{-q/p}, & k\leq N,\\
%         0, & k>N.
%     \end{cases}
% \]
% Then, for every $n\geq N$, we have
% \[
% \sum_{k=1}^{n}w_k a_k=V_N.
% \]
% Therefore,
% \[
% \sum_{n=1}^{\infty}u_n \!\left(\sum_{k=1}^{n}w_k a_k\right)^{\!p} = \sum_{n=1}^{N-1}u_n \!\left(\sum_{k=1}^{n}w_k a_k\right)^{\!p} +  V_N^p\sum_{n=N}^{\infty}u_n \geq V_N^p\sum_{n=N}^{\infty}u_n .
% \]
% Moreover, since $p(q-1)=q$, we have
% \[
% \sum_{k=1}^{\infty}v_k a_k^p = \sum_{k=1}^{N}v_k w_k^{p(q-1)}v_k^{-q} = \sum_{k=1}^{N}w_k^q v_k^{-q/p} = V_N .
% \]
% Using the assumed inequality, we obtain
% \[
% V_N^p\sum_{n=N}^{\infty}u_n \leq C^p V_N.
% \]
% Thus
% \[
% \left(\sum_{n=N}^{\infty}u_n\right)^{\!\!1/p} \!V_N^{1/q} \leq C < \infty.
% \]
% Taking the supremum over $N$ gives $\beta\leq C$. Hence
% \[
% \beta\leq C_{\mathrm{opt}} .
% \]

% We now prove sufficiency. Assume that $\beta<\infty$. Let
% $0<\theta<1$. For each $n\geq 1$, we write
% \[
% \sum_{k=1}^{n}w_k a_k = \sum_{k=1}^{n} \left(v_k^{1/p}a_k V_k^{\theta/q}\right) \!\left(w_k v_k^{-1/p}V_k^{-\theta/q}\right).
% \]
% By Hölder's inequality,
% \[
% \left(\sum_{k=1}^{n}w_k a_k\right)^{\!p} \leq \left( \sum_{k=1}^{n}v_k a_k^p V_k^{\theta(p-1)} \right)\! \left( \sum_{k=1}^{n}w_k^q v_k^{-q/p}V_k^{-\theta} \right)^{\!p/q}.
% \]
% By definition, $W_k=w_k^q v_k^{-q/p}$ and $V_k=\sum_{j=1}^{k}W_j$. Hence, $W_k = V_k-V_{k-1}$ and
% \begin{align*}
%     \sum_{k=1}^{n}W_k V_k^{-\theta} &= \sum_{k=1}^{n}(V_k-V_{k-1}) V_k^{-\theta} = \sum_{k=1}^{n}\int_{V_{k-1}}^{V_k}V_k^{-\theta}\,dt \\
%     %
%     &\leq  \sum_{k=1}^{n}\int_{V_{k-1}}^{V_k} t^{-\theta}\,dt = \int_0^{V_n} t^{-\theta}\,dt = \frac{V_n^{1-\theta}}{1-\theta}.
% \end{align*}
% Therefore,
% \[
% \left(\sum_{k=1}^{n}w_k a_k\right)^{\!p} \leq \frac{V_n^{(1-\theta)(p-1)}}{(1-\theta)^{p-1}} \sum_{k=1}^{n}v_k a_k^p V_k^{\theta(p-1)} .
% \]
% Multiplying by $u_n$ and summing over $n$, we get
% \begin{align*}
%     \sum_{n=1}^{\infty}u_n \!\left(\sum_{k=1}^{n}w_k a_k\right)^{\!p}
%     &\leq \frac{1}{(1-\theta)^{p-1}} \sum_{n=1}^{\infty} u_n V_n^{(1-\theta)(p-1)} \sum_{k=1}^{n}v_k a_k^p V_k^{\theta(p-1)} \\
%     %
%     &= \frac{1}{(1-\theta)^{p-1}} \sum_{k=1}^{\infty} v_k a_k^p V_k^{\theta(p-1)} \sum_{n=k}^{\infty} u_n V_n^{(1-\theta)(p-1)} .
% \end{align*}
% Let $\gamma:=(1-\theta)(p-1)$. We claim that
% \begin{equation}\label{eq - bound2}
%     \sum_{n=k}^{\infty}u_nV_n^{\gamma} \leq \frac{\beta^p}{\theta} V_k^{-\theta(p-1)} .
% \end{equation}
% Indeed, by summation by parts and by \eqref{eq - bound},
% \[
% \begin{aligned}
%     \sum_{n=k}^{\infty}u_nV_n^\gamma &= \sum_{n=k}^{\infty} (R_{n}-R_{n+1})V_n^\gamma = \sum_{n=k}^{\infty} R_{n}V_n^\gamma - \sum_{n=k}^{\infty} R_{n+1}V_n^\gamma \\
%     %
%     &= R_{k}V_k^\gamma + \sum_{n=k+1}^{\infty} R_{n}V_n^\gamma - \sum_{n=k+1}^{\infty} R_{n}V_{n-1}^\gamma \\
%     %
%     &\leq \beta^p V_k^{\gamma+1-p} + \beta^p \sum_{n=k+1}^{\infty} V_n^{1-p} \!\left(V_n^\gamma-V_{n-1}^\gamma\right).
% \end{aligned}
% \]
% Since $1-p<0$, the last sum is bounded above by the corresponding integral:
% \begin{align*}
%     \sum_{n=k+1}^{\infty} V_n^{1-p} \!\left(V_n^\gamma-V_{n-1}^\gamma\right) &= \sum_{n=k+1}^{\infty} \gamma\int_{V_{n-1}}^{V_n} V_n^{1-p}t^{\gamma-1} \,dt \leq \sum_{n=k+1}^{\infty} \gamma\int_{V_{n-1}}^{V_n} t^{1-p}t^{\gamma-1} \,dt \\
%     %
%     &= \gamma\int_{V_{k}}^{V_\infty} t^{\gamma-p} \,dt =\frac{\gamma}{p-1-\gamma}\bigl(V_k^{\gamma+1-p} - V_\infty^{\gamma+1-p} \bigr) \\
%     %
%     &\leq  \frac{\gamma}{p-1-\gamma} V_k^{\gamma+1-p},
% \end{align*}
% where we used the fact that $V_\infty:=\lim_{n\to\infty} V_n \geq 0$ and $p-1-\gamma=\theta(p-1)>0$.
% % \[
% %     \sum_{n=k+1}^{\infty}
% %     V_n^{1-p}
% %     \left(V_n^\gamma-V_{n-1}^\gamma\right)
% %     \leq
% %     \int_{V_k}^{\infty} t^{1-p}\,d(t^\gamma).
% % \]
% % Because $\gamma<p-1$, we have
% % \[
% %     \int_{V_k}^{\infty} t^{1-p}\,d(t^\gamma)
% %     =
% %     \gamma\int_{V_k}^{\infty}t^{\gamma-p}\,dt
% %     =
% %     \frac{\gamma}{p-1-\gamma}V_k^{\gamma+1-p}.
% % \]
% Hence
% \[
% \begin{aligned}
%     \sum_{n=k}^{\infty}u_nV_n^{\gamma}
%     &\leq
%     \beta^p\left(1+\frac{\gamma}{p-1-\gamma}\right) V_k^{\gamma+1-p} = \frac{p-1}{p-1-\gamma}\,\beta^p\,V_k^{\gamma+1-p}.
% \end{aligned}
% \]
% Since $\gamma=(1-\theta)(p-1)$, we have
% \[
% \frac{p-1}{p-1-\gamma} = \frac{1}{\theta} \quad\text{ and }\quad \gamma+1-p=-\theta(p-1),
% \]
% proving \eqref{eq - bound2}. Substituting \eqref{eq - bound2} into the previous estimate gives
% \[
% \begin{aligned}
%     \sum_{n=1}^{\infty}u_n\!\left(\sum_{k=1}^{n}w_k a_k\right)^{\!p} &\leq \frac{\beta^p}{\theta(1-\theta)^{p-1}} \sum_{k=1}^{\infty} v_k a_k^p V_k^{\theta(p-1)} V_k^{-\theta(p-1)} \\
%     %
%     &= \frac{\beta^p}{\theta(1-\theta)^{p-1}} \sum_{k=1}^{\infty}v_k a_k^p .
% \end{aligned}
% \]
% Optimizing over $0<\theta<1$, we find that the minimum is attained at $\theta=1/p$, where
% \[
% \frac{1}{\theta(1-\theta)^{p-1}} = p q^{p-1}.
% \]
% Consequently,
% \[
% \sum_{n=1}^{\infty}u_n \!\left(\sum_{k=1}^{n}w_k a_k\right)^{\!p} \leq p q^{p-1}\beta^p \sum_{k=1}^{\infty}v_k a_k^p .
% \]
% Taking $p$-th roots yields
% \[
% C_{\mathrm{opt}} \leq p^{1/p}q^{1/q}\beta .
% \]
% Together with the necessity proved above, this completes the proof.
% \end{proof}

Combining \Cref{prop - equiv,prop - old} immediately gives the following corollary.


\begin{cor}\label{cor - charac}
Let $1\leq p\leq q\leq\infty$. Let $(X,\mu)$ be a measure space, with $\mu$ semi-finite if $p=1$. Let $X = X_1 \cup X_2 \cup \cdots$ be a partition of $X$ and set $X^{(n)} := X_1 \cup \cdots \cup X_n$. Let $(b_n)_{n\geq 1}$ and $(M_n)_{n\geq 1}$ be sequences of weights. Let $m$ be a measurable function on $X$ and define $w_n$ as in \eqref{def - wn}. Then
\begin{equation*}\label{eq - charac-bound}
    \left(\sum_{n=1}^{\infty} b_n \left| \frac{1}{M_n} \int_{X^{(n)}}m(x)f(x)\,\mathrm{d}\mu(x) \right|^q \right)^{\!\!1/q}\!\! \leq p^{1/q}(p')^{1/p'} \beta \left( \sum_{n=1}^{\infty} b_n \int_{X_n}|f(x)|^p\,\mathrm{d}\mu(x) \right)^{\!\!1/p}
\end{equation*}
if and only if
\[
\beta = \sup_{N\ge1} \left(\sum_{n=N}^{\infty} \frac{b_n}{M_n^q}\right)^{\!\!1/q}\! \left(\sum_{k=1}^N \frac{w_k^{p'}}{b_k^{p'/p}}\right)^{\!\!1/p'}\!
<\infty
\]
\!$($\!with the obvious modifications if $p,q\in\{1,\infty\}$, as in \Cref{rem - explain}\!$)$.
\end{cor}

\begin{proof}
    By \Cref{prop - equiv}, the inequality is equivalent to the discrete inequality 
    \begin{equation}\label{eq - simple}
         \left(\sum_{n=1}^{\infty} \left|\frac{b_n^{1/q}}{M_n} \sum_{k=1}^{n}w_k a_k \right|^q \right)^{\!\!1/q}\! \leq \rho \left(\sum_{k=1}^{\infty} \big|b_k^{1/p}a_k\big|^p\right)^{\!\!1/p}.
    \end{equation}
    As we did in the proof of \Cref{prop - equiv}, observe that if $w_k= 0$, then the $k$th coordinate does not contribute to the left-hand side, and we may assume without any loss of generality that $a_k=0$. This is identical to just ignoring the $k$th summand in both the left and the right-hand sides, which is equivalent to simply assuming that $w_k\neq 0$. Hence, if we make this hypothesis, we can replace $a_k$ by $a_k/w_k$ and define 
    \[
    u_n=\frac{b_n^{1/q}}{M_n}
    \qquad\text{and}\qquad
    v_n=\frac{b_n^{1/p}}{w_n},
    \]
    to find that \eqref{eq - simple} is in fact equivalent to 
    \[
    \left(\sum_{n=1}^{\infty} \bigg| u_n\sum_{k=1}^n a_k \bigg|^q \right)^{\!\!1/q}\!\! \le \rho \left(\sum_{k=1}^{\infty}|v_ka_k|^p \right)^{\!\!1/p}.
    \]
    Hence, by \Cref{prop - old}, the desired inequality is equivalent to having
    \[
    \beta := \sup_{N\ge1} \left(\sum_{n=N}^{\infty} \frac{b_n}{M_n^q}\right)^{\!\!1/q}\! \left(\sum_{k=1}^N \frac{w_k^{p'}}{b_k^{p'/p}}\right)^{\!\!1/p'}\!
    <\infty .
    \]
    Moreover, we may choose $\rho = p^{1/q}(p')^{1/p'}\beta .$
\end{proof}

\begin{rem}
    In the important case where $b_n\equiv1$, for which the right-hand side of the inequality reduces to the $L^p(\mu)$-norm of $f$, the condition becomes simply
    \[
    \sup_{N\geq 1} \left( \sum_{n=N}^{\infty}\frac{1}{M_n^q} \right)^{\!\!1/q}\! \left(\sum_{k=1}^{N}w_k^{p'} \right)^{\!\!1/p'} \!< \infty.
    \]
\end{rem}



\Cref{cor - charac} characterizes the condition for which there is a $C>0$ such that
\[
\left(\sum_{n=1}^{\infty} b_n \left| \frac{1}{M_n} \int_{X^{(n)}}m(x)f(x)\,\mathrm{d}\mu(x) \right|^q \right)^{\!\!1/q}\!\! \leq C \left( \sum_{n=1}^{\infty} b_n \int_{X_n}|f(x)|^p\,\mathrm{d}\mu(x) \right)^{\!\!1/p}.
\]
However, this constant is not always optimal. Consider for instance the following cases.

\begin{enumerate}
    \item \emph{The classical discrete Hardy inequality}. Suppose that $X=\mathbb{N}$, $X_n=\{n\}$, $\mu$ is the counting measure, $p=q$, $m\equiv 1$, $b_n=1$ for all $n\geq 1$ and $M_n=n$. Then
    \[
    \beta = \sup_{N\geq 1} \bigg( \sum_{n=N}^{\infty}\frac{1}{n^p} \bigg)^{\!\!1/p}\! N^{1/q} = \sup_{N\geq 1} \bigg( N^{p-1}\sum_{n=N}^{\infty}\frac{1}{n^p} \bigg)^{\!\!1/p}\!.
    \]
    Now, define
    \[
    \varphi(x)=x\left(1-\Big(\frac{x}{x+1}\Big)^{p-1}\right).
    \]
    Then, with $y=(x+1)^{-1}$,
    \[
        \varphi'(x)=1-(1-y)^{p-1}(1+(p-1)y)\geq 0,
    \]
    since $(1-y)^{-a}\geq 1+ay$. Thus $\varphi$ is increasing, and for
    $n\geq N+1$,
    \[
    \frac{N((N+1)^{p-1}-N^{p-1})}{(N+1)^{p-1}} =\varphi(N)\leq \varphi(n) =n\!\left(1-\Big(\frac{n}{n+1}\Big)^{p-1}\right).
    \]
    Hence,
    \[
    \frac{(N+1)^{p-1}-N^{p-1}}{n^{p}} \leq \frac{(N+1)^{p-1}}{N} \left(\frac{1}{n^{p-1}}-\frac{1}{(n+1)^{p-1}}\right).
    \]
    Summing over $n\geq N+1$ gives
    \[
    ((N+1)^{p-1}-N^{p-1})\sum_{n=N+1}^{\infty}\frac{1}{n^{p}} \leq \frac{(N+1)^{p-1}}{N}\cdot\frac{1}{(N+1)^{p-1}} = \frac{1}{N}.
    \]
    Therefore, if $s_N:=N^{p-1}\sum_{n=N}^{\infty}\frac{1}{n^p}$, we find
    \[
    s_N-s_{N+1} = \frac{1}{N} - \bigl((N+1)^{p-1}-N^{p-1}\bigr) \sum_{n=N+1}^{\infty}\frac{1}{n^p} \geq 0,
    \]
    where $\zeta(\cdot)$ is the Riemann zeta function. Therefore, $s_N$ is decreasing and
    \[
    \beta =  \bigg(\sum_{n=1}^{\infty}\frac{1}{n^p} \bigg)^{\!\!1/p} \!= \zeta^{1/p}(p) < \infty.
    \]
    In particular, \Cref{cor - charac} yields 
    \[
    \sum_{n=1}^{\infty} \left( \frac{a_1+\cdots+a_n}{n} \right)^p \leq \left(\frac{p}{p-1}\right)^{\!p} (p-1)\zeta(p) \sum_{n=1}^{\infty} a_n^p,
    \]
    which is Hardy's inequality with the optimal constant multiplied by a factor of $(p-1)\zeta(p)$.

    \item \emph{The geometric Hardy inequality.} For every integer $b\geq 2$, \cite{Bouthat2023} shows that
    \[
    \sum_{k=1}^\infty \frac{1}{b^k} \bigg| \sum_{n=1}^{b^k} a_n \bigg|^2  \!\leq \frac{\sqrt{b}+1}{\sqrt{b}-1} \sum_{n=1}^\infty |a_n|^2 ,
    \]
    with the constant $\frac{\sqrt{b}+1}{\sqrt{b-1}} $ being optimal. With the corresponding choice of parameters, \Cref{cor - charac} gives the analogous inequality, but now with the constant 
    \[
    C = \frac{4b}{b-1}.
    \]
    Once again, the provided constant is not optimal.
\end{enumerate}


While \Cref{eq - charac-bound} is often not optimal, it is worth noting that it will never be far from optimal. Indeed, \Cref{prop - old} ensures that if $C_{\mathrm{opt}}$ denotes the optimal multiplicative constant for the inequality, then
\[
\beta \leq C_{\mathrm{opt}} \leq p^{1/q}(p')^{1/p'}\beta .
\]
Since for $1\leq p\leq q \leq \infty$, we have
\[
1 \leq p^{1/q}(p')^{1/p'} \leq p^{1/p}(p')^{1/p'} \leq 2,
\]
with the upper bound being attained only when $p=q=2$. Therefore, the bound of \Cref{cor - charac} can be off by at most a factor of 2.

% \subsection{Case $p = 1$}
% We have similar results when $p = 1$. The conjugate exponent to $p = 1$ is now $q = \infty$, which corresponds to the essential supremum norm.
% \begin{prop}
%     Let $X$ be a measurable space with a semi-finite measure $\mu$. Let $(b_n)_{n \geq 1}$ and $(M_n)_{n \geq 1}$ be sequences of positive numbers. Let $m$ be a positive measurable function on $X$ and $X = X_1 \cup X_2 \cup \cdots$ be a partition of $X$. % such that for every $k \in \mathbb{N}$ and every measurable subset $F \subseteq X_k$, there exists a measurable subset $E \subseteq F$ such that $0 < \mu (E) < \infty$. 
%     Set $X^{(n)} := X_1 \cup X_2 \cup \cdots \cup X_n$. Suppose that the quantities $w_n := \operatorname*{ess\,sup}_{x \in X_n} m (x)$ satisfy $w_n < \infty$ for any $n \geq 1$. For some $\rho \in (0, \infty )$, the inequality
%     \begin{equation}
%         \sum_{n = 1}^\infty b_n \left| \frac{1}{M_n} \int_{X^{(n)}} m (x) f (x) \,\mathrm{d}\mu (x) \right| \leq \rho \sum_{n = 1}^\infty b_n \int_{X_n} |f (x)| \,\mathrm{d}\mu (x) \label{eq:Pequals1EquivalenceFunctions}
%     \end{equation}
%     holds for every complex-valued measurable function $f$ on $X$ if and only if the discrete inequality
%     \begin{equation}
%         \sum_{n = 1}^\infty b_n \left| \frac{1}{M_n} \sum_{k = 1}^n w_k a_k \right| \leq \rho \sum_{k = 1}^\infty b_k |a_k| \label{eq:Pequals1EquivalenceSequences}
%     \end{equation}
%     holds for every sequence $(a_k)_{k \geq 1}$ of complex numbers.
% \end{prop}
% \begin{proof}
%     Assume that the inequality is true for sequences. It is sufficient to prove the inequality for positive and measurable functions because of the triangle inequality. Let $f$ be a positive measurable function defined on $X$. Set $a_k := \int_{X_k} f (x) \,\mathrm{d}\mu (x)$, for $k \in \mathbb{N}$. Then, since all the quantities are positive, we obtain
%         $$  
%             \sum_{n = 1}^\infty \frac{b_n}{M_n} \sum_{k = 1}^n w_k \int_{X_k} f (x) \,\mathrm{d}\mu (x)  \leq \rho \sum_{k = 1}^\infty b_k \int_{X_k} f (x) \,\mathrm{d}\mu (x) 
%         $$
%     and, after rearranging the left-hand side, we obtain \eqref{eq:Pequals1EquivalenceFunctions}.

%     Now suppose that the inequality is true for complex-valued measurable functions. From the triangle inequality, it suffices to prove the inequality for sequences of positive numbers. So assume that $(a_k)_{k \geq 1}$ is a sequence of positive numbers. We may also assume, without loss of generality, that $w_k > 0$. Let $\varepsilon > 0$. From the assumption on the semi-finiteness of the measure $\mu$ on $X_k$ and the definition of the essential supremum, there exists a measurable subset $E_k \subseteq X_k$ such that
%         \begin{enumerate}
%             \item $m (x) \geq w_k - \varepsilon$ for $x \in E_k$;
%             \item $\mu (E_k) > 0$. 
%         \end{enumerate}
%     Define $f_k := \frac{a_k}{\mu (E_k)} \chi_{E_k}$ for $k$ in the support of the sequence $(a_k)_{k \geq 1}$. Since all terms are positive, we have
%         $$  
%             \int_{X_k} m (x) f_k (x) \,\mathrm{d}\mu (x) = \int_{E_k} \frac{m (x) a_k}{\mu (E_k)} \,\mathrm{d}\mu (x) \geq (w_k - \varepsilon ) a_k 
%         $$
%     and therefore, for any $N \in \mathbb{N}$,
%         $$
%             \sum_{n = 1}^N b_n \left( \frac{1}{M_n} \sum_{k = 1}^n (w_k - \varepsilon ) a_k \right) \leq \sum_{n = 1}^N b_n \left( \frac{1}{M_n} \int_{\mathcal{X}_k} m (x) f (x) \,\mathrm{d}\mu (x) \right) 
%         $$
%     where $f := \sum_{k} f_k$. Since all the quantities are positive, taking the limit as $N \rightarrow \infty$, we deduce that
%     $$
%         \sum_{n = 1}^\infty b_n \left( \frac{1}{M_n} \sum_{k = 1}^n (w_k - \varepsilon ) a_k \right) \leq \sum_{n = 1}^\infty b_n \left( \frac{1}{M_n} \int_{X^{(n)}} m (x) f (x) \,\mathrm{d}\mu (x) \right) .
%     $$
%     The inequality for functions then implies that
%     $$
%          \sum_{n = 1}^\infty b_n \left( \frac{1}{M_n} \sum_{k = 1}^n (w_k - \varepsilon ) a_k \right) \leq \rho \sum_{k = 1}^\infty b_k \int_{X_k} f (x) \,\mathrm{d}\mu (x) = \rho \sum_{n = 1}^\infty b_k a_k .
%     $$
%     This last inequality is true for any $\varepsilon > 0$ and therefore we get the desired inequality for sequences. 
% \end{proof}

% By analyzing the quantities involved in the preceding theorem when $p = 1$ and as $q \rightarrow \infty$, we obtain the following extension to the case $p = 1$.  \textcolor{red}{[Probably move this before the examples]}
% \begin{prop}
%     Let $(u_n)_{n \geq 1}$, $(v_n)_{n \geq 1}$, and $(w_n)_{n \geq 1}$ be positive sequences. Set 
%         $$  
%             \beta_1 := \sup_{N \geq 1} \max_{1 \leq k \leq N} \frac{w_k}{v_k} \left( \sum_{n = N}^\infty u_n \right) .
%         $$
%     Then the inequality 
%         $$  
%             \sum_{n = 1}^\infty u_n \left| \sum_{k = 1}^n w_k a_k \right| \leq C \sum_{k = 1}^\infty v_k |a_k|
%         $$
%     holds for all sequences $a = (a_k)_{k \geq 1}$ and some $C > 0$ if and only if $\beta_1 < \infty$. Moreover, $C_{\mathrm{opt}} = \beta_1$. 
% \end{prop}
% \begin{proof}
%     Applying the triangle inequality to the inner sum of the left-hand side of the goal inequality, the result will follow by proving the equivalence for positive sequences. 

%     Suppose that $\beta_1 < \infty$. Then, for each $N \in \mathbb{N}$, the quantity $\max_{1 \leq k \leq N} \frac{w_k}{v_k} \sum_{n = N}^\infty u_n$ is finite. Fix $N \in \mathbb{N}$ for now. Using Tonelli's theorem, we have
%     $$  
%         \sum_{n = 1}^N u_n \sum_{k = 1}^n w_k a_k = \sum_{k = 1}^N w_k a_k \sum_{n = k}^N u_n .
%     $$
%     and therefore
%     \begin{align*}
%         \sum_{n = 1}^N u_n \sum_{k = 1}^n w_k a_k = \sum_{k=1}^N \left( \frac{w_k}{v_k} \sum_{n = k}^N u_n \right) v_k a_k &\leq \sum_{k = 1}^N \left( \max_{1 \leq j \leq k} \frac{w_j}{v_j} \sum_{n = k}^N u_n \right) v_k a_k \\ 
%         &\leq \sum_{k = 1}^N \left( \max_{1 \leq j \leq k} \frac{w_j}{v_j} \sum_{n = k}^\infty u_n \right) v_k a_k .
%     \end{align*}
%     Using the assumption, we deduce
%     $$
%         \sum_{n = 1}^N u_n \sum_{k = 1}^n w_k a_k \leq \beta_1 \sum_{k = 1}^N v_k a_k .
%     $$
%     Letting $N \rightarrow \infty$, we obtain the desired result with $C = \beta_1$. 

%     Now suppose that the inequality holds with a constant $C > 0$. For $m \in \mathbb{N}$, define the sequence $(a_k)_{k \geq 1}$ as follows
%         $$
%             a_k = \begin{cases}
%                 \big( \max_{1 \leq j \leq m} \frac{w_j}{v_j} \big) \tfrac{1}{w_m} & k = m \\ 
%                 0 & k \neq m .
%             \end{cases}
%         $$
%     Then, the inequality from the assumption applied to the sequence $(a_k)_{k \geq 1}$ implies that
%         $$  
%             w_m a_m \sum_{n = m}^\infty u_n \leq C v_m a_m 
%         $$
%     which can be rewritten as
%         %$$
%         %    \max_{1 \leq j \leq m} \frac{w_j}{v_j} \sum_{n = m}^\infty u_n \leq C \frac{v_m}{w_m} \max_{1 \leq j \leq m} \frac{w_j}{v_j} .
%         %$$
%     %Simplifying the last inequality, we deduce that
%         \begin{equation}
%             \frac{w_m}{v_m} \sum_{n = m}^\infty u_n \leq C . \label{eq:basicInequalityPEquals1}
%         \end{equation}
%     This last inequality holds for any $m \in \mathbb{N}$. Using \eqref{eq:basicInequalityPEquals1} for $1 \leq l \leq m$ in combination with the fact that $u_n > 0$, we obtain  
%         $$  
%             \frac{w_l}{v_l} \sum_{n = m}^\infty u_n \leq C
%         $$
%     and therefore
%         $$  
%             \max_{1 \leq l \leq m} \frac{w_l}{v_l} \sum_{n = m}^\infty u_n \leq C .
%         $$
%     Since $m$ was arbitrary, we finally have $\beta_1 < \infty$. This also proves the optimality of $\beta_1$.
% \end{proof}

% \begin{rem}
%     Notice that the proof of the last proposition suggests that the quantity $\beta_1$ may be replaced with an equivalent quantity, namely
%     $$
%         \beta_1' := \sup_{N \geq 1} \frac{w_N}{v_N} \sum_{n = N}^\infty u_n . 
%     $$
%     This last quantity is easier to handle in applications than the quantity in the statement of the proposition.
% \end{rem}

% As a corollary, we obtain an analogue result for the general case.
% \begin{cor}
%     Let $X$ be a measure space with a semi-finite measure $\mu$. Let $(b_n)_{n \geq 1}$ and $(M_n)_{n \geq 1}$ be sequences of positive numbers. Let $m$ be a positive measurable function defined on $X$. Let $X = X_1 \cup X_2 \cup \cdots$ be a partition of $X$ and set $X^{(n)} := X_1 \cup X_2 \cup \cdots \cup X_n$. Suppose that the quantity
%         $$
%             w_n = \operatorname*{ess\,sup}_{x \in X_n} m (x) < \infty .
%         $$
%     for every $n \geq 1$. Then the inequality
%         $$ 
%             \sum_{n = 1}^\infty b_n \left| \frac{1}{M_n} \int_{X^{(n)}} m (x) f (x) \,\mathrm{d}\mu (x) \right| \leq C \sum_{n = 1}^\infty b_n \int_{X_n} |f (x)| \,\mathrm{d}\mu (x)
%         $$
%     holds for some $C > 0$ and for every complex-valued measurable function $f$ on $X$ if and only if 
%         $$
%             \beta_1 := \sup_{N \geq 1} \frac{w_N}{b_N} \sum_{n = N}^\infty \frac{b_n}{M_n} < \infty .
%         $$
%     Moreover, we have $C_{\mathrm{opt}} = \beta_1$. 
% \end{cor}





% \textcolor{MidnightBlue}{INTRODUCTION: %(1) we introduce the problem while relating it to operator theory (2) we mention that it has been generalized (3) we mention that while it has been generalized, there is still room for apparent generalization and previous work has not given a characterization for boundedness (important because of operator theory) (4) we mention our proposition for generalization [maybe say that there are several ways to generalize, we focus on one and briefly explore another one, but we leave the full exploration as an open problem] (5) say that while it seems like more of a generalization, it can be reduced to the discrete case as it will be shown (6) give the outline of the paper 
% MAIN BODY: %(1) Give the equivalence between full case and discrete case (2) provide the full characterization and the corollary (3) show that while it characterizes, it does not always give a sharp constant with some examples (compare to previous work and to classical Hardy) (4) Finish the section and move to another section, where we look at a more local condition which is more restrictive be often give better bounds 
% (5) present what we proved last summer (6) present the examples (7) Finish the section and present some preliminary work on other generalizations on function spaces (8) conclusion and open questions}

% {\color{purple}
% \section{Idée de structure}

% On pourrait premièrement faire une intro assez classique qui motive l'intéret de partitionner l'espace en répétant très rapidement l'intro de mon autre article, et ensuite en disant qu'il y a des problèmes évidents avec l'autre approche puisque ça ne permet pas d'obtenir l'inégalité classique de Hardy.

% \smallskip
% La première section pourrait ensuite se concentrer sur le cas discret et être essentiellement ce qu'on a en ce moment comme résultats.

% \smallskip
% Ensuite, on pourrait faire une autre section qui généralise le plus possible (tout en restant pratique et élégant) les résultats de la Section 1. Il y aurait entre autre le cas avec l'intégrale (qu'on pourrait redémontrer au complet, ou bien juste expliquer comment adapter l'argument du Théorème 1), potentiellement un autre cas avec les deux intégrales, et des généralisations que j'ai vu en ligne que je suis sûr qu'on peut adapter. 

% \smallskip
% On pourrait ensuite conclure en présentant tous les cas particulier qu'on peut trouver dans la littérature qu'on couvre avec notre approche, et en présentant quelques questions ouvertes.

% \smallskip
% (Évidemment, toutes ces questions peuvent être approfondies en regardant l'optimalité ou non des inégalités).
% }

\section{Local sufficient conditions}\label{sec - local}

As described above, \Cref{cor - charac} provides a complete characterization of the cases for which the inequality in \Cref{eq - charac-bound} holds. While we showed that the provided bound has a multiplicative constant that cannot be off by a factor of more than $2$ (and more precisely $p^{1/q}(p')^{1/p'}$), one may seek to find another approach for which the bounds obtained are often sharper, at least in the more important cases. 

Moreover, the condition in \Cref{cor - charac} is not always easy to verify. In this section, we provide another approach to treat the main inequality of this paper. This approach, based on Broadbent’s proof of Hardy's inequality \cite{MR1574000}, gives a simple sufficient condition to verify while also having the benefit of being sharp in several important cases.

\begin{thm}\label{thm - main}
    Let $1< p\leq q<\infty$. Let $(X,\mu)$ be a measure space. Let $X = X_1 \cup X_2 \cup \cdots$ be a partition of $X$ and set $X^{(n)} := X_1 \cup \cdots \cup X_n$. Let $(b_n)_{n\geq 1}$ and $(M_n)_{n\geq 1}$ be sequences of weights. Let $m$ be a measurable function on $X$ and define $w_n$ as in \eqref{def - wn}. Suppose that $w_n>0$ and set $\gamma_n:=b_n^{1/p+1/q'}\!\!/w_n$. If $0<\rho<\infty$ is such that
    \[
    \frac{1}{q}\frac{M_n(\gamma_n-\gamma_{n+1})}{b_n}
    +\frac{1}{q'}\frac{\gamma_n(M_n-M_{n-1})}{b_n}
    \geq
    \frac{1}{\rho}
    \qquad \text{for }~ n=1,2,3,\dots,
    \]
    then
    \begin{align*}
        \left(
        \sum_{n = 1}^\infty b_n
        \left| \frac{1}{M_n} \int_{X^{(n)}} m (x) f (x)
        \,\mathrm{d}\mu (x) \right|^q
        \right)^{\!\!1/q}
        \leq
        \rho
        \left(
        \sum_{n = 1}^\infty b_n \int_{X_n} |f (x)|^p
        \,\mathrm{d}\mu (x)
        \right)^{\!\!1/p}.
    \end{align*}
\end{thm}

\begin{proof}
    First observe that we always have
    \[
    \sum_{n = 1} b_n \left| \frac{1}{M_n} \int_{X^{(n)}} m (x) f (x) \mathrm{d}\mu (x) \right|^q \leq \sum_{n = 1}^\infty b_n \!\left( \frac{1}{M_n} \int_{X^{(n)}} |m (x) f(x)|  \mathrm{d}\mu (x) \right)^{\!q}.
    \]
    Hence, it is enough to prove the estimate with $m$ and $m$ replaced by $|m|$ and $|f|$, since $w_n$ is defined in terms of $|m|$ and the right-hand side depends only on $|f|$.
    %
    Define $M_0=T_0:=0$, let $T_n := \frac{1}{M_n} \int_{X^{(n)}} |m(x) f (x)| \,\mathrm{d}\mu (x)$ for $n\geq 1$, and let
    \[
    \Delta_n := b_n T_n^q - \rho\gamma_n w_n\left(\int_{X_n} |f(x)|^p \,\mathrm{d}\mu(x)\right)^{\!\!1/p} T_n^{q-1}.
    \]
    Then Hölder's inequality implies that
    \begin{align*}
        \int_{X_n} |m(x) f(x)| \,\mathrm{d}\mu(x)
        &\leq \left(\int_{X_n} |m(x)|^{p'} \,\mathrm{d}\mu(x) \right)^{\!\!1/p'}\!
        \left(\int_{X_n} |f(x)|^p \,\mathrm{d}\mu(x) \right)^{\!\!1/p} \\
        %
        &= w_n \left(\int_{X_n} |f(x)|^p \,\mathrm{d}\mu(x) \right)^{\!\!1/p}.
    \end{align*}
    It thus follows that
    \begin{align*}
        \Delta_n &= b_n T_n^q - \rho\gamma_nw_n\left(\int_{X_n} |f(x)|^p \,\mathrm{d}\mu(x)\right)^{\!\!1/p} T_n^{q-1} \\
        %
        &\leq  b_n T_n^q - \rho\gamma_n \left( \int_{X_n} |m(x) f(x)| \,\mathrm{d}\mu(x) \right) T_n^{q-1} \\
        %
        &= b_n T_n^q - \rho\gamma_n \big(M_nT_n-M_{n-1} T_{n-1}\big) T_n^{q-1} \\
        %
        &= T_n^q \left( b_n - \rho\gamma_n M_n \right)
        + \rho\gamma_n M_{n-1} T_{n-1} T_{n}^{q-1}.
    \end{align*}
    But now recall Young's inequality which implies that
    \[
    T_{n-1} T_{n}^{q-1} \leq \frac{1}{q}\, T_{n-1}^q +  \frac{1}{q'}\, T_n^q.
    \]
    Hence, it follows that
    \begin{align*}
        \Delta_n &\leq  T_n^q \left( b_n - \rho\gamma_n M_n \right)
        + \rho\gamma_n M_{n-1} T_{n-1} T_{n}^{q-1} \\
        %
        &\leq T_n^q \left( b_n - \rho\gamma_n M_n
        + \frac{q-1}{q}\rho\gamma_n M_{n-1} \right)
        + \frac{\rho\gamma_n M_{n-1}}{q} T_{n-1}^q.
    \end{align*}
    Now, observe that
    \begin{gather*}
        b_n - \rho\gamma_n M_n
        + \frac{q-1}{q}\rho\gamma_n M_{n-1}
        \leq -\frac{\rho\gamma_{n+1} M_n}{q},
        \qquad \forall n\geq 1 \\
        %
        \Updownarrow \\
        %
        \frac{1}{q}\frac{M_n(\gamma_n-\gamma_{n+1})}{b_n}
    +\frac{1}{q'}\frac{\gamma_n(M_n-M_{n-1})}{b_n}
        \geq
        \frac{1}{\rho},
        \qquad \forall n\geq 1 ,
    \end{gather*}
    which is satisfied by hypothesis. Hence, we have
    \begin{align*}
        \Delta_n &\leq T_n^q \left( b_n - \rho\gamma_n M_n
        + \frac{q-1}{q}\rho\gamma_n M_{n-1} \right)
        + \frac{\rho\gamma_n M_{n-1}}{q} T_{n-1}^q \\
        %
        &\leq \frac{\rho\gamma_n M_{n-1}}{q} T_{n-1}^q
        - \frac{\rho\gamma_{n+1} M_n}{q} T_n^q.
    \end{align*}
    Now, this forms a telescopic series. Hence, since $M_0=0$, it follows that
    \begin{align*}
        \sum_{n=1}^s \Delta_n
        &\leq \sum_{n=1}^s
        \left( \frac{\rho\gamma_n M_{n-1}}{q} T_{n-1}^q
        - \frac{\rho\gamma_{n+1} M_n}{q} T_n^q \right) %\\
        %
        %&= \frac{\rho\gamma_1 M_0}{q}T_0^q
        %- \frac{\rho\gamma_{s+1}M_s}{q}T_s^q
        =
        - \frac{\rho\gamma_{s+1}M_s}{q}T_s^q \leq 0.
    \end{align*}
    Hence, since
    \[
        \Delta_n
        =
        b_nT_n^q
        -
        \rho\gamma_nw_n
        \left(\int_{X_n} |f(x)|^p \,\mathrm{d}\mu(x)\right)^{\!\!1/p}
        T_n^{q-1},
    \]
    we have by letting $s\to\infty$ that
    \[
    \sum_{n=1}^\infty b_n T_n^q
    \leq
    \rho \sum_{n=1}^\infty \gamma_n w_n
    \left( \int_{X_n} |f(x)|^p \,\mathrm{d}\mu(x) \right)^{\!\!1/p}
    T_n^{q-1}.
    \]
    Therefore, another application of Hölder's inequality finally implies that
    \begin{align*}
        \sum_{n=1}^\infty b_n T_n^q
        &\leq \rho \sum_{n=1}^\infty \gamma_nw_n
        \!\left(\int_{X_n} |f(x)|^p \,\mathrm{d}\mu(x) \right)^{\!\!1/p}
        T_n^{q-1} \\
        %
        &=
        \rho \sum_{n=1}^\infty
        \left[
        b_n^{1/p}
        \left(\int_{X_n} |f(x)|^p \,\mathrm{d}\mu(x) \right)^{\!\!1/p}
        \right]
        \left[
        b_n^{1/q'}T_n^{q-1}
        \right] \\
        %
        &\leq \rho
        \left(\sum_{n=1}^\infty b_n \int_{X_n} |f(x)|^p \,\mathrm{d}\mu(x) \right)^{\!\!1/p}\!
        \left(
        \sum_{n=1}^\infty
        b_n^{p'/q'}T_n^{(q-1)p'}
        \right)^{\!\!1/p'}.
    \end{align*}
    Since $p\leq q$, we have $\frac{q-1}{q}\geq \frac{1}{p'}.$ Therefore,
    \begin{align*}
        \left(
        \sum_{n=1}^\infty
        b_n^{p'/q'}T_n^{(q-1)p'}
        \right)^{\!\!1/p'}\!
        =
        \left(
        \sum_{n=1}^\infty
        (b_nT_n^q)^{\frac{(q-1)p'}{q}}
        \right)^{\!\!1/p'}\! \leq
        \left(
        \sum_{n=1}^\infty b_nT_n^q
        \right)^{\!\!1-1/q}.
    \end{align*}
    Hence,
    \begin{align*}
        \sum_{n=1}^\infty b_n T_n^q
        &\leq
        \rho
        \left(\sum_{n=1}^\infty b_n \int_{X_n} |f(x)|^p \,\mathrm{d}\mu(x) \right)^{\!\!1/p}\!
        \left(
        \sum_{n=1}^\infty b_nT_n^q
        \right)^{\!\!1-1/q},
    \end{align*}
    from which it finally follows that
    \begin{align*}
         \left( \sum_{n=1}^\infty b_n T_n^q \right)^{\!\!1/q}\!
         \leq
         \rho
         \left(\sum_{n=1}^\infty b_n \int_{X_n} |f(x)|^p \,\mathrm{d}\mu(x) \right)^{\!\!1/p},
    \end{align*}
    as desired.
\end{proof}




% \begin{thm}\label{thm - main}
%     Let $1<p<\infty$. Let $X$ be a measure space with measure $\mu$. Let $X = X_1 \cup X_2 \cup \cdots$ be a partition of $X$ and set $X^{(n)} := X_1 \cup \cdots \cup X_n$. Let $(b_n)_{n\geq 1}$ and $(M_n)_{n\geq 1}$ be sequences of positive numbers. Let $m$ be a positive measurable function defined on $X$ and $f$ a complex measurable function on $X$, and suppose that 
%     \begin{equation}\label{eq - def_wn}
%         w_n := \Big(\int_{X_n} m(x)^q \,\mathrm{d}\mu (x) \Big)^{1/q}.
%     \end{equation}
%     Suppose that $0<\rho<\infty$ is such that  
%     \[
%     \left(1-\frac{w_nb_{n+1}}{b_nw_{n+1}}\right)\!M_n+(p-1)(M_n-M_{n-1}) \geq \frac{pw_n}{\rho} \qquad \text{for }~ n=1,2,3,\dots
%     \]
%     Then
%     \begin{align*}
%         \sum_{n = 1}^\infty b_n \left| \frac{1}{M_n} \int_{X^{(n)}} m (x) f (x)  \,\mathrm{d}\mu (x) \right|^p \leq \rho^p \sum_{n = 1}^\infty b_n \int_{X_n} |f (x)|^p  \,\mathrm{d}\mu (x)
%     \end{align*}
% \end{thm}
% \begin{proof}
%     First observe that we always have
%     \[
%     \sum_{n = 1} b_n \left| \frac{1}{M_n} \int_{X^{(n)}} m (x) f (x) \mathrm{d}\mu (x) \right|^p \leq \sum_{n = 1}^\infty b_n \!\left( \frac{1}{M_n} \int_{X^{(n)}} m (x) |f(x)|  \mathrm{d}\mu (x) \right)^{\!p}.
%     \]
%     Hence, without loss of generality, we may assume that $f(x)\geq0$ for all $x\in X$. 
%     %
%     Define $M_0=T_0:=0$, let $T_n := \frac{1}{M_n} \int_{X^{(n)}} m(x) f (x) \,\mathrm{d}\mu (x)$ for $n\geq 1$, and let 
%     \[
%     \Delta_n := b_n T_n^p - \rho  b_n\left(\int_{X_n} f^p(x) \,\mathrm{d}\mu(x)\right)^{\!\!1/p} T_n^{p-1}.
%     \]
%     Then Hölder's inequality implies that
%     \begin{align*}
%         \int_{X_n} m(x) f(x) \,\mathrm{d}\mu(x) &\leq \left(\int_{X_n} m^q(x) \,\mathrm{d}\mu(x) \right)^{\!\!1/q} \left(\int_{X_n} f^p(x) \,\mathrm{d}\mu(x) \right)^{\!\!1/p} \\
%         %
%         &= w_n \left(\int_{X_n} f^p(x) \,\mathrm{d}\mu(x) \right)^{\!\!1/p},
%     \end{align*}
%     where the $w_n$ are defined as in \eqref{eq - def_wn}. It thus follows that
%     \begin{align*}
%         \Delta_n &= b_n T_n^p - \rho  b_n\left(\int_{X_n} f^p(x) \,\mathrm{d}\mu(x)\right)^{\!\!1/p} T_n^{p-1} \\
%         %
%         &\leq  b_n T_n^p - \frac{\rho b_n}{w_n} \left( \int_{X_n} m(x) f(x) \,\mathrm{d}\mu(x) \right) T_n^{p-1} \\
%         %
%         &= b_n T_n^p - \frac{\rho b_n}{w_n} \big(M_nT_n-M_{n-1} T_{n-1}\big) T_n^{p-1} \\
%         %
%         &= T_n^p \left( b_n - \frac{\rho b_n M_n}{w_n} \right) + \frac{\rho b_n M_{n-1}}{w_n} T_{n-1} T_{n}^{p-1}.
%     \end{align*}
%     But now recall Young's inequality which implies that
%     \[
%     T_{n-1} T_{n}^{p-1} \leq \frac{T_{n-1}^p}{p} + (p-1) \frac{T_n^p}{p}
%     \]
%     for all $0<t<\infty$. Hence, it follows that
%     \begin{align*}
%         \Delta_n &\leq  T_n^p \left( b_n - \frac{\rho b_n M_n}{w_n} \right) + \frac{\rho b_n M_{n-1}}{w_n} T_{n-1} T_{n}^{p-1} \\
%         %
%         &\leq T_n^p \left( b_n - \frac{\rho b_n M_n}{w_n} + \frac{p-1}{p}\frac{\rho b_n M_{n-1}}{w_n} \right) + \frac{\rho b_n  M_{n-1}}{pw_n} T_{n-1}^p.
%     \end{align*}
%     Now, observe that
%     \begin{gather*}
%         b_n - \frac{\rho b_n M_n}{w_n} + \frac{p-1}{p}\frac{\rho b_n M_{n-1}}{w_n}   \leq -\frac{\rho b_{n+1}  M_{n}}{pw_{n+1}}, \qquad \forall n\geq 1 \\
%         %
%         \Updownarrow \\
%         %
%         \left(1-\frac{w_nb_{n+1}}{b_nw_{n+1}}\right)\!M_n+(p-1)(M_n-M_{n-1}) \geq \frac{pw_n}{\rho} , \qquad \forall n\geq 1 ,
%     \end{gather*}
%     which is satisfied by hypothesis. Hence, we have
%     \begin{align*}
%         \Delta_n &\leq T_n^p \left( b_n - \frac{\rho b_n M_n}{w_n} + \frac{p-1}{p}\frac{\rho b_n M_{n-1}}{w_n} \right) + \frac{\rho b_n  M_{n-1}}{pw_n} T_{n-1}^p \\
%         %
%         &\leq \frac{\rho b_n  M_{n-1}}{pw_n} T_{n-1}^p - \frac{\rho b_{n+1}  M_{n}}{pw_{n+1}} T_{n}^p.
%     \end{align*}
%     Now, this forms a telescopic series and it follows that
%     \begin{align*}
%         \sum_{n=1}^s \Delta_n &\leq \sum_{n=1}^s \left( \frac{\rho b_n  M_{n-1}}{pw_n} T_{n-1}^p - \frac{\rho b_{n+1}  M_{n}}{pw_{n+1}} T_{n}^p \right) \\
%         %
%         &= \frac{\rho b_1  M_{0}}{pw_1} T_{0}^p - \frac{\rho b_{s+1}  M_{s}}{pw_{s+1}} T_{s}^p = - \frac{\rho b_{s+1}  M_{s}}{pw_{s+1}} T_{s}^p \\
%         %
%         &\leq 0.
%     \end{align*}
%     Hence, since $\Delta_n = b_n T_n^p - \rho  b_n\left(\int_{X_n} f^p(x) \,\mathrm{d}\mu(x)\right)^{\!\frac{1}{p}} T_n^{p-1}$, we have by letting $s\to\infty$ that
%     \[
%     \sum_{n=1}^\infty b_n T_n^p \leq \rho \sum_{n=1}^\infty b_n \left( \int_{X_n} f^p(x) \,\mathrm{d}\mu(x) \right)^{\!\!1/p} T_n^{p-1}.
%     \]
%     Therefore, another application of Hölder's inequality finally implies that
%     \begin{align*}
%         \sum_{n=1}^\infty b_n T_n^p &\leq \rho \sum_{n=1}^\infty  b_n \!\left(\int_{X_n} f^p(x) \,\mathrm{d}\mu(x) \right)^{\!\frac{1}{p}} T_n^{p-1} \\
%         %
%         &\leq \rho \left(\sum_{n=1}^\infty b_n \int_{X_n} f^p(x) \,\mathrm{d}\mu(x) \right)^{\!\frac{1}{p}} \!\left( \sum_{n=1}^\infty b_n T_n^p \right)^{\!\frac{1}{q}} 
%     \end{align*}
%     from which it finally follows that 
%     \begin{align*}
%          \left( \sum_{n=1}^\infty b_n T_n^p \right)^{\!\frac{1}{p}} \leq \rho \left(\sum_{n=1}^\infty b_n \int_{X_n} f^p(x) \,\mathrm{d}\mu(x) \right)^{\!\frac{1}{p}},
%     \end{align*}
%     as desired.
% \end{proof}



\begin{rem}\label{rem - counter-example}
    The condition of \Cref{thm - main} is not necessary. Indeed, consider the case $X=\mathbb{N}$, $X_n=\{n\}$, $\mu$ the discrete measure, $p=q$, $b_n\equiv 1$, $M_n=\alpha^{n}$ and $m_n=\beta^{-n}$ with $\alpha,\beta>1$ and $\frac{\beta}{q}+\frac{1}{\alpha q'}=1$. Then
    \begin{align*}
        \frac{1}{q}\frac{M_n(\gamma_n-\gamma_{n+1})}{b_n}
    +\frac{1}{q'}\frac{\gamma_n(M_n-M_{n-1})}{b_n} = 0,
    \end{align*}
    which implies that $\rho=\infty$. However, 
    \begin{align*}
        \beta &\leq \sup_{N\geq 1} \left( \sum_{n=1}^{\infty}\frac{1}{\alpha^{np}} \right)^{\!\!1/p}\!\! \left(\sum_{k=1}^{\infty} \frac{1}{\beta^{kp'}}
    \right)^{\!\!1/p'}\! = \frac{1}{(\alpha^p-1)^{1/p}(\beta^{p'}-1)^{1/p'}} < \infty.
    \end{align*}
    Hence, \Cref{cor - charac} allows us to conclude that 
    \begin{align*}
        \left(\sum_{n=1}^{\infty} \Bigg|\frac{1}{M_n}\sum_{k \in \mathbf{N}_n} m_k a_k \Bigg|^p\right)^{\!\!1/p}
        \!\!\leq  \frac{p^{1/p}(p')^{1/p'}}{(\alpha^p-1)^{1/p}(\beta^{p'}-1)^{1/p'}}\!\left(\sum_{n=1}^\infty|a_n|^p\right)^{\!\!1/p}\!,
    \end{align*}
    while \Cref{thm - main} does not.
\end{rem}




Now, as we did previously, let us consider some important particular cases for which we recover some results in the literature.

\begin{enumerate}
    % \item \emph{The classical discrete Hardy inequality.} Suppose that $X=\mathbb{N}$, $X_n=\{n\}$, $\mu$ is the counting measure, $m\equiv 1$, $b_n=1$ for all $n\geq 1$ and $M_n=n$. Then $w_n = 1$ for all $n\geq 1$ and the sufficient condition becomes
    % \[
    % \rho \geq \frac{p}{p-1} \qquad \text{for }~ n=1,2,3,\dots
    % \]
    % Hence, choosing $\rho=\frac{p}{p-1}$ recover the sharp constant in the classical discrete Hardy inequality.

    \item \emph{The general discrete Hardy inequality.} Choose $X=\mathbb{N}$, $X_n=\{n\}$, $\mu$ is the counting measure, $p=q$, $m_n=b_n$ for all $n\geq 1$ and $M_n=m_1+\cdots+m_n$. Then $w_n = b_n$ for all $n\geq 1$ and the sufficient condition once again becomes
    \[
    \rho \geq \frac{p}{p-1} \qquad \text{for }~ n=1,2,3,\dots
    \]
    Hence, choosing $\rho=\frac{p}{p-1}$ recovers the sharp constant in the weighted discrete Hardy inequality, due to Copson \cite{Copson1928}. The discrete Hardy inequality corresponds to the case $m_k\equiv 1$. Note that the results from both \cite{Bouthat2023} and \cite{VincentSohani2025} did not have the classical Hardy inequality as a special case. %Moreover, note that the realization of Copson's inequality is possible because of the addition of the weights $b_n$.

    \item \emph{The geometric Hardy inequality.} Let $p=q=2$, $X=\mathbb{N}$, $\mu$ the discrete measure, $X_1=\{1,\dots,b\}$ and $X_n=\{b^{n-1}+1,\dots,b^{n}\}$ for $n\geq 2$, where $b\geq 2$ is an integer. Moreover, let $b_n\equiv 1$, $m_n=1$ and $M_n = b^{n/2}$. Then, if $n\geq 2$, we have %we find that $w_n = \sqrt{b^n-b^{n-1}}$ if $n>1$ and $w_1 = \sqrt{b}$. Therefore, for $n\geq 2$, we have
    \begin{align*}
        \frac{1}{q}\frac{M_n(\gamma_n-\gamma_{n+1})}{b_n}
    +\frac{1}{q'}\frac{\gamma_n(M_n-M_{n-1})}{b_n} = \frac{\sqrt{b}-1}{\sqrt{b-1}}
    \end{align*}
    and
    \[
    \frac{1}{q}\frac{M_n(\gamma_n-\gamma_{n+1})}{b_n}
    +\frac{1}{q'}\frac{\gamma_n(M_n-M_{n-1})}{b_n} = 1-\frac{1}{2\sqrt{b-1}} \geq \frac{\sqrt{b}-1}{\sqrt{b-1}}
    \]
    if $n=1$. Hence, the optimal value for $\rho^2$ is $\bigl(\frac{\sqrt{b}-1}{\sqrt{b-1}}\bigr)^{-2} = \frac{\sqrt{b}+1}{\sqrt{b}-1}$, which yields
    \begin{equation*}\label{eq - expo}
         \sum_{k=1}^{\infty} \frac{1}{b^k}  \bigg| \sum_{j=1}^{b^k} a_j \bigg|^2
        \leq~
        \frac{\sqrt{b}+1}{\sqrt{b}-1} \sum_{n=1}^{\infty} |a_n|^2.
    \end{equation*}
    This is precisely \cite[Theorem 2]{Bouthat2023}, where it is shown that this inequality is sharp.

    \item \emph{Example 1 in \cite{Bouthat2023}}. Consider the case $p=q=2$, $X=\mathbb{N}$, $\mu$ the discrete measure, $b_n\equiv1$, $X_n=\{2^{n-1},\dots,2^n-1\}$, $M_n = \sum_{k=1}^n k\cdot 2^{(k-1)/2}$ and $m_n$ is the sequence which takes the value $n$ for all indices in $X_n$. Then $w_k=k\cdot 2^{(k-1)/2}$, and it is shown in \cite{Bouthat2023} that the bound holds with the constant $\rho \approx 1.2114$. However, 
    \begin{align*}
        \frac{1}{q}\frac{M_n(\gamma_n-\gamma_{n+1})}{b_n}
    +\frac{1}{q'}\frac{\gamma_n(M_n-M_{n-1})}{b_n}
        %
        &= \frac{1}{2} \!\bigg(1+2^{\frac{n}{2}}\bigg(\frac{1}{n}-\frac{1}{\sqrt{2}(n+1)}\bigg)\sum_{k=1}^{n}k\,2^{\frac{k}{2}} \bigg) \\
        %
        &\geq \frac{1}{2} \!\bigg(1+2^{\frac{1}{2}}\bigg(\frac{1}{1}-\frac{1}{\sqrt{2}(1+1)}\bigg)\sum_{k=1}^{1}k\,2^{\frac{k}{2}} \bigg) \\
        %
        &= \frac{6-\sqrt{2}}{4}.
    \end{align*}
    Hence, \Cref{thm - main} ensures that we can take the constant $\rho = \frac{4}{6-\sqrt{2}} \approx 0.87226$, which is a significant improvement on \cite[Example 1]{Bouthat2023}.
\end{enumerate}





% \begin{rem}\label{rem - Hardy}
%     The case where $N_n=\{n\}$ and $M_n = w_1+\cdots+w_n$ in the Corollary above correspond precisely to the classical weighted Hardy inequality, due to Copson in 1928 \cite{Copson1928}. Furthermore, if $m_k\equiv 1$, then we also get back the classical Hardy inequality. Note that \cite[Theorem 1]{Bouthat2023} did not have the classical Hardy inequality as a special case, while ours does.
% \end{rem}







% In the following, we follow the notation established in \ref{} and we define 
% \begin{equation}\label{eq - def_wn}
%     w_n:=\left(\sum_{k\in N_n} m_k^q \right)^{\!\frac{1}{q}}
% \end{equation}

% \begin{thm}\label{thm - main}
%     Let $(a_n)_{n \geq 1}$ be a sequence of complex numbers, let $(m_n)_{n \geq 1}$ and $(M_n)_{n \geq 1}$ be two sequences of weights, and let $w_n$ be defined as in \eqref{eq - def_wn}. Let $\mathbb{N} = N_1 \cup N_2 \cup \cdots$ be a partition of the natural numbers, and let $\mathbf{N}_n := N_1  \cup \cdots \cup N_n$. Let $p>1$ and let $q=\frac{p}{p-1}$ and suppose that
%         \[
%         \rho \,:=\, \inf_{n \geq 1}  \left\{\frac{M_{n}}{w_{n}}-\left(\frac{M_{n}}{w_{n+1}}\right)^{\!\frac{1}{p}}\!\!\left(\frac{M_{n-1}}{w_{n}}\right)^{\!\frac{1}{q}}\right\}^{\!-1} 
%         \]
%         is finite. Then
%         \begin{align*}
%         	\sum_{n=1}^{\infty} \Bigg|\frac{1}{M_n}\sum_{k \in \mathbf{N}_n} m_k a_k \Bigg|^p
%         	\leq~ \rho^p  \sum_{n=1}^{\infty} |a_n|^p.
%         \end{align*}
% \end{thm}
% \begin{proof}
%     Let $T_n := \frac{1}{M_n}\sum_{k \in \mathbf{N}_n}m_k a_k$, and let 
%     \[
%     \Delta_n := T_n^p - \rho \left(\sum_{k\in N_n} |a_k|^p \right)^{\frac{1}{p}} T_n^{p-1}.
%     \]
%     Then Hölder's inequality implies that
%     \[
%     \sum_{k\in N_n} m_k|a_k| \leq \left(\sum_{k\in N_n} m_k^q \right)^{\!\frac{1}{q}} \left(\sum_{k\in N_n} |a_k|^p \right)^{\!\frac{1}{p}} = w_n \left(\sum_{n\in N_n} |a_n|^p \right)^{\!\frac{1}{p}},
%     \]
%     where the $w_n$ are defined as in \eqref{eq - def_wn}, and it follows that
%     \begin{align*}
%         \Delta_n :\!&= T_n^p - \rho \left(\sum_{k\in N_n} |a_k|^p \right)^{\frac{1}{p}} T_n^{p-1} \leq  T_n^p - \frac{\rho}{w_n} \left(\sum_{k\in N_n} |a_k| \right) T_n^{p-1} \\
%         %
%         &= T_n^p - \frac{\rho}{w_n} \big(M_nT_n-M_{n-1} T_{n-1}\big) T_n^{p-1} \\
%         %
%         &= T_n^p \left( 1 - \frac{\rho M_n}{w_n} \right) + \frac{\rho M_{n-1}}{w_n} T_{n-1} T_{n}^{p-1}.
%     \end{align*}
%     But now recall Young's inequality which implies that for all $0<t_n<\infty$,
%     \[
%     T_{n-1} T_{n}^{p-1} \leq \frac{t_n^pT_{n-1}^p}{p} + (p-1) \frac{T_n^p}{pt_n^q}.
%     \]
%     Hence, it follows that
%     \begin{align*}
%         \Delta_n &\leq  T_n^p \left( 1 - \frac{\rho M_n}{w_n} \right) + \frac{\rho M_{n-1}}{w_n} T_{n-1} T_{n}^{p-1} \\
%         %
%         &\leq T_n^p \left( 1 - \frac{\rho M_n}{w_n} + \frac{p-1}{p}\frac{\rho M_{n-1}}{w_n t_n^q} \right) + \frac{\rho t_n^p M_{n-1}}{pw_n} T_{n-1}^p.
%     \end{align*}
%     Choose $t_n:=\left(\frac{M_{n-1}w_{n+1}}{M_{n}w_{n}}\right)^{\frac{1}{p+q}}$ and observe that 
%     \begin{gather*}
%         1 - \frac{\rho M_n}{w_n} + \frac{p-1}{p}\frac{\rho M_{n-1}}{w_n t_n^q} \leq -\frac{\rho t_n^p M_{n}}{pw_{n+1}}, \qquad \forall n\geq 1 \\
%         %
%         \Updownarrow \\
%         %
%         \rho \geq \frac{w_n}{M_{n}\!\left(1-\frac{w_n t_n^p}{p w_{n+1}}\right)-\frac{M_{n-1}}{q t_n^q}} , \qquad \forall n\geq 1 \\
%         %
%         \Updownarrow \\
%         %
%         \rho \geq \frac{1}{\frac{M_{n}}{w_{n}}\!\left(1-\left(\frac{M_{n-1}}{w_{n}}\frac{w_{n+1}}{M_{n}}\right)^{\frac{1}{q}}\right)}, \qquad \forall n\geq 1 ,
%     \end{gather*}
%     which is trivially satisfied by hypothesis. Hence, we have
%     \begin{align*}
%         \Delta_n &\leq T_n^p \left( 1 - \frac{\rho M_n}{w_n} + \frac{p-1}{p}\frac{\rho M_{n-1}}{w_n t_n^q} \right) + \frac{\rho t_n^pM_{n-1}}{pw_n} T_{n-1}^p \\
%         %
%         &\leq \frac{\rho t_n^p M_{n-1}}{pw_n} T_{n-1}^p - \frac{\rho t_n^p M_{n}}{pw_{n+1}} T_{n}^p.
%     \end{align*}
%     Now, this forms a telescopic series and it follows that
%     \begin{align*}
%         \sum_{n=1}^s \Delta_n &\leq \sum_{n=1}^s \left( \frac{\rho t_n^p M_{n-1}}{pw_n} T_{n-1}^p - \frac{\rho t_n^p M_{n}}{pw_{n+1}} T_{n}^p \right) \\
%         %
%         &= \frac{\rho t_0^p M_{0}}{pw_1} T_{0}^p - \frac{\rho t_s^p M_{s}}{pw_{s+1}} T_{s}^p \\
%         %
%         &= - \frac{\rho t_s^p M_{s}}{pw_{s+1}} T_{s}^p \\
%         %
%         &\leq 0.
%     \end{align*}
%     Therefore, since $\Delta_n = T_n^p - \rho \left(\sum_{n\in N_n} |a_n|^p \right)^{\frac{1}{p}} T_n^{p-1}$, we finally have by letting $s\to\infty$ that
%     \[
%     \sum_{n=1}^\infty T_n^p \leq \rho \sum_{n=1}^\infty \left(\sum_{k\in N_n} |a_k|^p \right)^{\!\frac{1}{p}} T_n^{p-1}.
%     \]
%     Hence, another application of Hölder's inequality implies that
%     \begin{align*}
%         \sum_{n=1}^\infty T_n^p &\leq \rho \sum_{n=1}^\infty \left(\sum_{k\in N_n} |a_k|^p \right)^{\!\frac{1}{p}} T_n^{p-1} \\
%         %
%         &\leq \rho \left(\sum_{n=1}^\infty \sum_{k\in N_n} |a_k|^p \right)^{\!\frac{1}{p}} \left( \sum_{n=1}^\infty T_n^p \right)^{\!\frac{1}{q}} \\
%         %
%         &= \rho \left(\sum_{n=1}^\infty |a_n|^p \right)^{\!\frac{1}{p}} \left( \sum_{n=1}^\infty T_n^p \right)^{\!\frac{1}{q}}
%     \end{align*}
%     from which it finally follows that 
%     \begin{align*}
%          \left( \sum_{n=1}^\infty T_n^p \right)^{\!\frac{1}{p}} \leq \rho \left(\sum_{n=1}^\infty |a_n|^p \right)^{\!\frac{1}{p}},
%     \end{align*}
%     which concludes the proof.
% \end{proof}







%%%%%%%%%%%%%%%%%%%%%%%%%%%%%%%%%%%%%
%%%%%%% Calcul Hardy classique %%%%%%
%%%%%%%%%%%%%%%%%%%%%%%%%%%%%%%%%%%%%

% \begin{example}
%     Now, observe that if we let $m_n=1$, $N_n:=\{n\}$ and $M_n:=n$, then we have $w_n=1$ and
% \begin{align*}
%     \rho &= \sup_{n \geq 1}  \left(\frac{w_n}{M_{n}\!\left(1-\frac{w_n}{p w_{n+1}}\right)-\frac{M_{n-1}}{q}}\right) \\
%     %
%     &= \sup_{n \geq 1}  \left(\frac{1}{n\!\left(1-\frac{1}{p}\right)-\frac{n-1}{q}}\right) \\
%     %
%     &= \sup_{n \geq 1}  \,q \\
%     %
%     &= q.
% \end{align*}
% % \begin{align*}
% %     \rho = \inf_{n \geq 1}  \left\{\frac{M_{n}}{w_{n}}-\left(\frac{M_{n}}{w_{n+1}}\right)^{\!\frac{1}{p}}\!\!\left(\frac{M_{n-1}}{w_{n}}\right)^{\!\frac{1}{q}}\right\}^{\!-1} &= \inf_{n \geq 1}  \left\{n-n^{\frac{1}{p}}(n-1)^{\frac{1}{q}}\right\}^{\!-1} \\
% %     %
% %     &= \left\{ \lim_{n\to\infty} \left(n-n^{\frac{1}{p}}(n-1)^{\frac{1}{q}}\right)\right\}^{\!-1} \\
% %     %
% %     &= q.
% % \end{align*}
% Hence, in this special case we retrieve the classical Hardy inequality.
% \end{example}



% \begin{example}\label{example 2}
% Let $p=2$, $b_n\equiv 1$, $m_n=1$, $M_n = b^{n/2}$ and $N_1=\{1,\dots,b\}$, $N_n=\{b^{n-1}+1,\dots,b^{n}\}$ for $n\geq 2$, where $b\geq 2$ is an integer, then we find that $w_n = \sqrt{b^n-b^{n-1}}$ if $n>1$ and $w_1 = \sqrt{b}$. Therefore, for $n\geq 2$, we have
% \begin{align*}
%     \frac{1}{pw_n}&\left(\left(1-\frac{w_nb_{n+1}}{b_nw_{n+1}}\right)\!M_n+(p-1)(M_n-M_{n-1})\right) = \frac{\sqrt{b-1}}{\sqrt{b}+1},
% \end{align*}
% while
% \[
% \frac{1}{pw_n}\left(\left(1-\frac{w_nb_{n+1}}{b_nw_{n+1}}\right)\!M_n+(p-1)(M_n-M_{n-1})\right) = 1-\frac{1}{2\sqrt{b-1}} \geq \frac{\sqrt{b-1}}{\sqrt{b}+1}.
% \]
% if $n=1$. Therefore, we find in that case that 
% \begin{equation}\label{eq - expo}
%      \sum_{k=1}^{\infty} \frac{1}{b^k}  \Bigg| \sum_{j=1}^{b^k} a_j \Bigg|^2
%     \leq~
%     \frac{\sqrt{b}+1}{\sqrt{b-1}} \sum_{n=1}^{\infty} |a_n|^2,
% \end{equation}
% which is precisely Theorem 2 in our last paper, where we proved that this inequality is sharp. Hence, this new inequality seems to somewhat bridge the gap between the classical Hardy inequality and our inequality. We also note that while \eqref{eq - expo} was proven in \cite{Bouthat2023}, it was not as a direct consequence of the main theorem. Indeed, it required the same techniques but with some different and more careful estimation. On the other hand, we obtain this elegant inequality directly as a consequence of our \Cref{thm - main}.
% \end{example}




% \begin{example}
%     Consider Example 1 in \cite{Bouthat2023}, namely the case $b_n\equiv1$, $p=2$, $N_n=\{2^{n-1},\dots,2^n-1\}$, $M_n = \sum_{k=1}^n k\cdot 2^{(k-1)/2}$ and $m_n$ is the sequence which takes the value $n$ for all indices in $N_n$. Then $w_k=k\cdot 2^{(k-1)/2}$ and it is shown in \cite{Bouthat2023} that $\rho \approx 1.2114$. However, in our case we find that
%     \begin{align*}
%         \frac{1}{pw_n}\!&\left(\left(1-\frac{w_nb_{n+1}}{b_nw_{n+1}}\right)\!M_n+(p-1)(M_n-M_{n-1})\right) \\
%         %
%         &= \frac{1}{2} \!\left(1+2^{\frac{n}{2}}\!\bigg(\frac{1}{n}-\frac{1}{\sqrt{2}(n+1)}\bigg)\sum_{k=1}^{n}k\,2^{\frac{k}{2}} \right) \\
%         %
%         &\geq \frac{1}{2} \!\left(1+2^{\frac{1}{2}}\!\bigg(\frac{1}{1}-\frac{1}{\sqrt{2}(1+1)}\bigg)\sum_{k=1}^{1}k\,2^{\frac{k}{2}} \right) \\
%         %
%         &= \frac{6-\sqrt{2}}{4} \approx 1.1464.
%     \end{align*}
%     Hence, we can take $\rho = \frac{4}{6-\sqrt{2}} \approx 0.87226$, which is better than what was obtained previously.
%     % \begin{align*}
%     %     \rho &= \sup_{n\geq 1} \left( \frac{2^{\frac{n-1}{2}}n}{\frac{1}{2}\left(\sum_{k=1}^{n}2^{\frac{k-1}{2}}k\left(1-\frac{n}{\sqrt{2}\left(n+1\right)}\right)+2^{\frac{n-1}{2}}n\right)}\right) \\
%     %     %
%     %     &= \sup_{n\geq 1} \left( \frac{1}{1+\frac{\left(4+3\sqrt{2}\right)\left(2^{-n/2}-1\right)+(1+\sqrt{2})n2^{-n/2}}{2n(1+n)}}\right) \\
%     %     %
%     %     &= \frac{1}{1+\frac{\left(4+3\sqrt{2}\right)\left(2^{-1/2}-1\right)+(1+\sqrt{2})2^{-1/2}}{2(1+1)}} \\
%     %     %
%     %     &= \frac{4\left(8+\sqrt{2}\right)}{31} \\
%     %     %
%     %     &\approx 1.2147.
%     % \end{align*}
%     % While the two constant are incredibly close, the one obtained in \cite{} is indeed smaller than our constant.
% \end{example}

% The above examples might make one think that the inequality of \Cref{thm - main} is always better than the one of \cite[Theorem 1]{Bouthat2023}. However, this is not the case, as the following example shows.

% \textcolor{red}{[Rephrase to show that is is not necessary]}

% \begin{example}\label{ex - counter-example}
%     Consider the case $X=\mathbb{N}$, $\mu$ the discrete measure, $b_n\equiv 1$, $M_n:=p^{n}$, $m_n=\beta^{-n}$, where $\beta=p-\frac{1}{q}\in(1,\infty)$, and $X_n=\{n\}$. Then $w_n=m_n=\beta^{-n}$ and we find that 
%     \begin{align*}
%         \left(1-\frac{w_nb_{n+1}}{b_nw_{n+1}}\right)\!M_n+(p-1)(M_n-M_{n-1}) = 0,
%     \end{align*}
%     which implies that $\rho=\infty$. However, 
%     \[
%     \beta = \sup_{N\geq 1} \left( \sum_{n=N}^{\infty}\frac{1}{p^{np}} \right)^{\!\!1/p}\! \left(\sum_{k=1}^{N} \beta^{-kq} 
%     \right)^{\!\!1/q} \!=  \sup_{N\geq 1} \frac{p^{1 - N}}{(p^p-1)^{1/p}} \left(\frac{ 1 - \beta^{-N q}}{\beta^q-1} \right)^{\!\!1/q} \leq \frac{1}{(p^p-1)^{1/p}(\beta^q-1)^{1/q}} < \infty.
%     \]
%     Hence, \Cref{cor - charac} allows us to conclude that 
%     \begin{align*}
%         \sum_{n=1}^{\infty} \Bigg|\frac{1}{M_n}\sum_{k \in \mathbf{N}_n} m_k a_k \Bigg|^p
%         \leq~  \frac{pq^{p/q}}{(p^p-1)(\beta^q-1)^{p/q}} \sum_{n=1}^\infty|a_n|^p,
%     \end{align*}
%     while \Cref{thm - main} does not.
% \end{example}

The above examples might make one think that the inequality of \Cref{thm - main} is always better than the one of \cite[Theorem 1]{Bouthat2023}. However, this is not the case, as the following example shows.

\begin{example}
    Consider the case $X=\mathbb{N}$, $\mu$ the discrete measure, $X_n=\{n\}$, $p=q$, $b_n\equiv 1$, $m_n=2^n$ and $M_n=\sum_{k=1}^n w_k = 2(2^n-1)$ with $1<p\leq \frac{5}{4}$. Then $w_n=m_n=2^n$ and \Cref{thm - main} yields 
    \begin{align*}
        \frac{1}{q}\frac{M_n(\gamma_n-\gamma_{n+1})}{b_n}
    +\frac{1}{q'}\frac{\gamma_n(M_n-M_{n-1})}{b_n}  = 1-\frac{1}{p2^{n}},
    \end{align*}
    and we can take $\rho = \frac{1}{1-1/2p}$. On the other hand, the constant given by \cite[Theorem 1]{Bouthat2023} (see \eqref{E:3}) is
    \begin{align*}
        \sup_{n\geq 1} \left(w_n\sum_{k=n}^\infty \frac{1}{M_n} \right)^{\!\!1/p}\!\! = \sup_{n\geq 1} \left(\sum_{k=0}^\infty \frac{1}{2(2^{k}-2^{-n})} \right)^{\!\!1/p}\!\! = \left(\sum_{k=1}^\infty \frac{1}{2^{k}-1}\right)^{\!\!1/p}\! \lessapprox 1.606695.
    \end{align*}
    Now, since $p\leq \frac{5}{4}$, we have $\rho = \frac{1}{1-1/2p} \geq \frac{1}{1-1/(2\cdot \frac{5}{4})} = \frac{5}{3} > 1.606695$. Hence, in this case \cite[Theorem 1]{Bouthat2023} gives a sharper constant than \Cref{thm - main}.
\end{example}


Let us make one last remark to conclude this section. Recall that \Cref{rem - counter-example} showed that there are some cases for which \cite[Theorem 1]{Bouthat2023} is able to provide an inequality (albeit not always with a sharp constant) while \Cref{thm - main} is not. However, it turns out that such a phenomenon only happens in some specific cases. Indeed, it turns out that in one of the more natural case, namely the case $X=\mathbb{N}$, $\mu$ the discrete measure, $p=q$, $b_n\equiv 1$ and $M_n=w_1+\cdots + w_n$, which is the subject of \cite[Theorem 1]{Bouthat2023}, our result subsumes the latter whenever it applies. 

More precisely, in this case, the sufficient condition of \Cref{thm - main} simplifies to
\begin{equation}\label{eq - reduced}
    \frac{1}{p}\left(\frac{1}{w_n}-\frac{1}{w_{n+1}}\right)\!M_n+\frac{1}{p'} \geq \frac{1}{\rho} \qquad \text{for }~ n=1,2,3,\dots
\end{equation}
Now, we claim that in this case, if
\[
\frac{1}{p}\left(\frac{1}{w_n}-\frac{1}{w_{n+1}}\right)\!M_n+\frac{1}{p'} > 0 \qquad \text{for }~ n=1,2,3,\dots,
\]
and if the estimate of the constant in \Cref{thm - main} diverges, then the estimate of \cite[Theorem 1]{Bouthat2023} is infinite. In other words, in this situation, the global condition of \cite[Theorem 1]{Bouthat2023} induces the local restriction of \Cref{thm - main}.


% {\color{red}
% In particular, if $b_n\equiv1$, then the inequality becomes
%         \begin{align*}
%         	\sum_{n = 1} \left| \frac{1}{M_n} \int_{X^{(n)}} m (x) f (x)  \mathrm{d}\mu (x) \right|^p \leq \rho^p \int_{X} |f (x)|^p  \mathrm{d}\mu (x)
%         \end{align*}
%         Moreover, if $M_n=w_1+\cdots+w_n$, then the condition becomes
%         \begin{equation}\label{eq - reduced}
%             \frac{1}p\left(\frac{1}{w_n}-\frac{1}{w_{n+1}}\right)\!M_n+\frac{1}{q} \geq \frac{1}{\rho} \qquad \text{for }~ n=1,2,3,\dots
%         \end{equation}
% }

% Nonetheless, it turns out that in the arguably most important and natural case, namely the case $b_n\equiv 1$ and $M_n=w_1+\cdots + w_n$, which is the subject of \cite[Theorem 1]{Bouthat2023}, our result subsume the latter whenever it applies. More precisely, on one hand, if
% \[
% \frac{1}p\left(\frac{1}{w_n}-\frac{1}{w_{n+1}}\right)\!M_n+\frac{1}{q} > 0 \qquad \text{for }~ n=1,2,3,\dots,
% \]
% and if the estimate of \cite[Theorem 1]{Bouthat2023} diverges, then our estimate of $\rho$ is infinite. On the other hand, the converse does not hold, as \Cref{rem - Hardy} shows. This is the content of the following proposition.

\begin{prop}\label{prop - infty}
    Consider the hypothesis of \Cref{thm - main} with $X=\mathbb{N}$, $\mu$ the discrete measure, $p=q$, $b_n\equiv 1$ and $M_n=w_1+\dots+w_n$. Moreover, suppose that 
    \begin{equation}
        \frac{1}{p}\left(\frac{1}{w_n}-\frac{1}{w_{n+1}}\right)\!M_n+\frac{1}{p'} > 0 \qquad \text{for }~ n=1,2,3,\dots %\tag*{($\dagger$)}
    \end{equation}
    If $\rho=\infty$, then 
    \[
    \rho' = \sup_{n\geq 1} \left( w_n \sum_{j=n}^{\infty} \frac{1}{M_j} \right) = \infty.
    \]
\end{prop}
\begin{proof}
    Since we suppose that 
    \[
    \frac{1}{p}\left(\frac{1}{w_n}-\frac{1}{w_{n+1}}\right)\!M_n+\frac{1}{p'} > 0 \qquad \text{for }~ n=1,2,3,\dots,
    \]
    and since $\rho$ is such that 
    \[
    \frac{1}{p}\left(\frac{1}{w_n}-\frac{1}{w_{n+1}}\right)\!M_n+\frac{1}{p'} \geq \frac{1}{\rho} \qquad \text{for }~ n=1,2,3,\dots,
    \]
    then 
    \begin{align}\label{eq - liminf}
        \rho=\infty \quad\iff\quad  \,\,&\frac{1}{p'} + \frac{1}{p} \liminf_{n\to\infty} \left(\frac{1}{w_n}-\frac{1}{w_{n+1}}\right)\!M_n =0 \nonumber\\
        %
        \iff\quad  \,\,&\liminf_{n\to\infty} \left(\frac{1}{w_{n+1}}-\frac{1}{w_n}\right)\!M_n = p-1.
    \end{align} 
    Denote by $\ell\in [0,\infty]$ the quantity $\ell:= \limsup_{n\to\infty}w_{n}$. If $\ell=\infty$, then $w_{n+1} > w_n$ infinitely often, which would imply that $\frac{1}{w_{n+1}}-\frac{1}{w_n} < 0$ infinitely often and contradict \eqref{eq - liminf}. Therefore, we must have $\ell<\infty$. 
    Since $\limsup_{n\to\infty} w_n = \ell$, there exists an $N_\varepsilon\in \mathbb{N}$ such that $w_n < \ell+\frac{\varepsilon}{2}$ for all $n\geq N_\varepsilon$. Hence, we have for all $n$ sufficiently large (without loss of generality, suppose that $n\geq N_\varepsilon$ suffice) that 
    \begin{align*}
        M_n &= w_1+\cdots +w_{N_\varepsilon-1} + w_{N_\varepsilon}  \cdots + w_n <  K_\varepsilon + (n-N_\varepsilon) \Bigl(\ell+\frac{\varepsilon}{2}\Bigr) \leq (\ell+\varepsilon) n
    \end{align*}
    Therefore, it finally follows that
    \begin{align*}
        \sup_{n\geq 1} \left( w_n \sum_{j=n}^{\infty} \frac{1}{M_j} \right) &\geq w_{N_\varepsilon} \sum_{j=N_\varepsilon}^{\infty} \frac{1}{M_j} > \frac{w_{N_\varepsilon}}{\ell+\varepsilon}  \sum_{j=N_\varepsilon}^{\infty} \frac{1}{j} = \infty,
    \end{align*} 
    as desired.
\end{proof}

% \begin{rem}
%     \Cref{prop - infty} cannot be extended to the most general case presented in \cite{Bouthat2023}, namely the case where we only have $b_n\equiv 1$. Indeed, \Cref{ex - counter-example} shows an example where \Cref{thm - main} yield $\rho = \infty$ while \cite[Theorem 1]{Bouthat2023} yields $\rho' < \infty$.
% \end{rem}



% \textcolor{red}{[Mention somewhere the difference between the local and global structure of the conditions in the new vs old theorem. This is something that our theorem is really bad at and which does not make sense to me and we should find another proof to solve this problem...]}



% {\color{red}
% The following corollary is the special case $b_n=w_n$.

% \begin{cor}
%     Let $(a_n)_{n \geq 1}$ be a sequence of complex numbers, let $(m_n)_{n \geq 1}$ and $(M_n)_{n \geq 1}$ be sequences of weights, and let $w_n$ be defined as in \eqref{eq - def_wn}. Let $\mathbb{N} = N_1 \cup N_2 \cup \cdots$ be a partition of the natural numbers, and let $\mathbf{N}_n := N_1  \cup \cdots \cup N_n$. Let $p>1$ and $q=\frac{p}{p-1}$, and suppose that $0<\rho<\infty$ is such that
%         \[
%         \frac{M_n-M_{n-1}}{w_n} \geq \frac{q}{\rho} \qquad\text{for }~ n=1,2,3,\dots
%         \]
%         Then
%         \begin{align*}
%         	\sum_{n=1}^{\infty} w_n \Bigg|\frac{1}{M_n}\sum_{k \in \mathbf{N}_n} m_k a_k \Bigg|^p
%         	\leq~ \rho^p \sum_{n=1}^\infty w_n \sum_{k\in N_n} |a_k|^p .
%         \end{align*}
%         In particular, if $M_n=w_1+\dots+w_n$, then 
%         \begin{align*}
%         	\sum_{n=1}^{\infty} w_n \Bigg|\frac{1}{M_n}\sum_{k \in \mathbf{N}_n} m_k a_k \Bigg|^p
%         	\leq~ q^p \sum_{n=1}^\infty w_n \sum_{k\in N_n} |a_k|^p,
%         \end{align*}
%         \textcolor{red}{and the constant $q^p$ is optimal.}
% \end{cor}
% \begin{proof}
%     The only nontrivial fact to show is the optimality of the last condition.
% \end{proof}
% }




% \section{A complete characterization}

% \begin{prop}
% Let $X$ be a measure space with measure $\mu$. Let $(b_n)_{n\geq 1}$ and $(M_n)_{n\geq 1}$ be sequences of positive numbers. Let $m$ be a positive measurable function defined on $X$ and $f$ a complex measurable function on $X$, and let $w_n$ be defined as in \eqref{eq - def_wn}. Let $X = X_1 \cup X_2 \cup \cdots$ be a partition of $X$ and set $X^{(n)} := X_1 \cup \cdots \cup X_n$. Let $1<p<\infty$ and suppose that 
% \[
%     w_n
%     :=
%     \left(
%         \int_{X_n}m(x)^q\,\mathrm{d}\mu(x)
%     \right)^{1/q}
%     <\infty
% \]
% for every $n\geq 1$. Let $(b_n)_{n\geq 1}$ and $(M_n)_{n\geq 1}$ be positive
% sequences. Then the inequality
% \begin{equation}
%     \sum_{n=1}^{\infty} b_n
%     \left|
%         \frac{1}{M_n}
%         \int_{X^{(n)}}m(x)f(x)\,\mathrm{d}\mu(x)
%     \right|^p
%     \leq
%     \rho^p
%     \sum_{n=1}^{\infty} b_n
%     \int_{X_n}|f(x)|^p\,\mathrm{d}\mu(x)
%     \label{eq - measure-hardy}
% \end{equation}
% holds for every measurable function $f$ if and only if the discrete inequality
% \begin{equation}
%     \sum_{n=1}^{\infty} b_n
%     \left|
%         \frac{1}{M_n}
%         \sum_{k=1}^{n}w_k a_k
%     \right|^p
%     \leq
%     \rho^p
%     \sum_{k=1}^{\infty} b_k |a_k|^p
%     \label{eq - discrete-hardy}
% \end{equation}
% holds for every sequence $a=(a_k)_{k\geq 1}$.
% % Consequently, the optimal constant in \eqref{eq - measure-hardy} is exactly
% % \[
% %     \rho_{\mathrm{opt}}
% %     =
% %     \sup_{a\neq 0}
% %     \frac{
% %         \left(
% %             \sum_{n=1}^{\infty} b_n
% %             \left|
% %                 \frac{1}{M_n}
% %                 \sum_{k=1}^{n}w_k a_k
% %             \right|^p
% %         \right)^{1/p}
% %     }{
% %         \left(
% %             \sum_{k=1}^{\infty} b_k |a_k|^p
% %         \right)^{1/p}
% %     }.
% % \]
% % Equivalently, $\rho_{\mathrm{opt}}$ is the norm on $\ell^p(b)$ of the weighted
% % Hardy matrix
% % \[
% %     A_{n,k}:=
% %     \frac{w_k}{M_n}\mathbf 1_{\{k\leq n\}} .
% % \]
% \end{prop}

% \begin{proof}
% We first prove that \eqref{eq - discrete-hardy} implies
% \eqref{eq - measure-hardy}. For each $k\geq 1$, define
% \[
%     F_k:=
%     \left(
%         \int_{X_k}|f(x)|^p\,\mathrm{d}\mu(x)
%     \right)^{1/p}.
% \]
% By Hölder's inequality,
% \[
%     \left|
%         \int_{X_k}m(x)f(x)\,\mathrm{d}\mu(x)
%     \right|
%     \leq
%     \left(
%         \int_{X_k}m(x)^q\,\mathrm{d}\mu(x)
%     \right)^{1/q}
%     \left(
%         \int_{X_k}|f(x)|^p\,\mathrm{d}\mu(x)
%     \right)^{1/p}
%     =
%     w_k F_k .
% \]
% Therefore,
% \[
%     \left|
%         \int_{X^{(n)}}m(x)f(x)\,\mathrm{d}\mu(x)
%     \right|
%     =
%     \left|
%         \sum_{k=1}^{n}
%         \int_{X_k}m(x)f(x)\,\mathrm{d}\mu(x)
%     \right|
%     \leq
%     \sum_{k=1}^{n}w_k F_k .
% \]
% Applying \eqref{eq - discrete-hardy} to the sequence $(F_k)_{k\geq 1}$ gives
% \[
% \begin{aligned}
%     \sum_{n=1}^{\infty} b_n
%     \left|
%         \frac{1}{M_n}
%         \int_{X^{(n)}}m(x)f(x)\,\mathrm{d}\mu(x)
%     \right|^p
%     &\leq
%     \sum_{n=1}^{\infty} b_n
%     \left|
%         \frac{1}{M_n}
%         \sum_{k=1}^{n}w_k F_k
%     \right|^p                                      \\
%     &\leq
%     \rho^p
%     \sum_{k=1}^{\infty}b_k F_k^p                  \\
%     &=
%     \rho^p
%     \sum_{k=1}^{\infty}b_k
%     \int_{X_k}|f(x)|^p\,\mathrm{d}\mu(x).
% \end{aligned}
% \]
% This proves \eqref{eq - measure-hardy}.

% Conversely, suppose that \eqref{eq - measure-hardy} holds. We prove \eqref{eq - discrete-hardy}. If $w_k=0$, then the $k$th coordinate does not contribute to the left-hand side of \eqref{eq - discrete-hardy}; hence we may first assume $a_k=0$ whenever $w_k=0$. For each $k$ such that $w_k>0$, define $f$ on $X_k$ by
% \[
% f(x) := a_k w_k^{1-q}m(x)^{q-1},
% \qquad x\in X_k .
% \]
% Then
% \[
% \int_{X_k}m(x)f(x)\,\mathrm{d}\mu(x) = a_kw_k^{1-q}
% \int_{X_k}m(x)^q\,\mathrm{d}\mu(x) = a_kw_k .
% \]
% Moreover, since $p(q-1)=q$, we have
% \[
% \begin{aligned}
%     \int_{X_k}|f(x)|^p\,\mathrm{d}\mu(x)
%     &=
%     |a_k|^p w_k^{p(1-q)}
%     \int_{X_k}m(x)^{p(q-1)}\,\mathrm{d}\mu(x)        \\
%     &=
%     |a_k|^p w_k^{p(1-q)}
%     \int_{X_k}m(x)^q\,\mathrm{d}\mu(x)              \\
%     &=
%     |a_k|^p w_k^{p(1-q)}w_k^q    \\
%     &=
%     |a_k|^p .
% \end{aligned}
% \]
% Applying \eqref{eq - measure-hardy} to this function $f$, we obtain
% \[
% \sum_{n=1}^{\infty} b_n
% \left|
%     \frac{1}{M_n}
%     \sum_{k=1}^{n}w_k a_k
% \right|^p
% \leq
% \rho^p
% \sum_{k=1}^{\infty} b_k |a_k|^p .
% \]
% This is precisely \eqref{eq - discrete-hardy}. The equality of the optimal
% constants follows immediately from the equivalence.
% \end{proof}









\section{Concluding remarks}\label{sec - conclusion}


In this paper, we generalized and characterized a class of Hardy-type inequalities in which the usual arithmetic means are replaced by weighted averages over nested subsets of a measure space. We did so using two different approaches, both distinct from the one previously used in the literature. 
%Our first contribution was to show that the measure-theoretic problem is, in a precise sense, equivalent to a discrete weighted Hardy inequality. This reduction shows that the boundedness of the corresponding operator depends only on the sequence of $L^q$-norms associated with the pieces of the partition, together with the prescribed weights. This equivalence allowed us to obtain a complete characterization of boundedness. More precisely, we proved that the inequality holds for some finite constant if and only if a natural testing quantity $B$ is finite. Moreover, if $C_{\mathrm{opt}}$ denotes the best possible constant, then \[ B \leq C_{\mathrm{opt}} \leq p^{1/p}q^{1/q}B. \] Consequently, the constant obtained from this characterization is always within a factor $p^{1/p}q^{1/q}\leq 2$ of the optimal one. This gives a general and robust answer to the boundedness question, although the resulting constant is not always sharp. To obtain sharper constants in important cases, we also developed a second method based on a local sufficient condition. This approach, inspired by Broadbent's proof of Hardy's inequality, is less general in the sense that its condition is not necessary. However, when it applies, it can give substantially better estimates. In particular, it recovers the sharp constant in the classical weighted Hardy inequality and in the geometric Hardy inequality, and it improves the constants obtained from previous general results in some examples. At the same time, 
Moreover, using some examples, we showed that the local criterion does not dominate the earlier global conditions in every possible situation. Thus the two approaches are complementary: the characterization gives a complete boundedness criterion, while the local condition often gives sharper constants in structured cases. 

Several natural questions remain open. First, it would be interesting to determine whether the present framework can be extended beyond $L^p$ spaces to more general function spaces, such as weighted $L^p$ spaces, Lorentz spaces, Orlicz spaces, or other Banach function spaces. Since the proof of the reduction relies heavily on Hölder's inequality and duality between $L^p$ and $L^{p'}$, such extensions would likely require identifying the correct analogue of the quantities $w_n$. 

A second direction concerns kernel operators. The inequalities considered here can be viewed as boundedness results for triangular averaging kernels. This suggests studying more general operators of the form 
\[ 
Tf(x)=\int_X k(x,y)f(y)\,\mathrm{d}\mu(y) 
\] 
and asking for conditions under which $T$ is bounded on $L^p$. Some cases may be accessible through Schur-type tests, but the classical Hardy operator already shows that Schur's test does not capture the full picture. It would therefore be useful to develop criteria adapted to Hardy-type kernels, especially those arising from nested sets or monotone integration regions. 

Finally, one could also consider versions in which the averages are replaced by sums of integrals over several pieces of a partition, or even over a sequence of different measure spaces. Such formulations would include inequalities of the form 
\[ 
\left(\sum_{n=1}^{\infty} b_n \bigg| \frac{1}{M_n} \sum_{k=1}^{n} \int_{X_k} m_k(x)f_k(x)\,\mathrm{d}\mu_k(x) \bigg|^q\right)^{\!\!1/q}\! \leq \rho \left( \sum_{n=1}^{\infty} b_n \int_{X_n}|f_n(x)|^p\,\mathrm{d}\mu_n(x) \right)^{\!\!1/p}. 
\] 
The methods of this paper suggest that such inequalities should again be closely related to discrete weighted Hardy inequalities, with the quantities $w_k$ replaced by suitable $L^{p'}$-norms of the functions $m_k$. Making this precise would provide a flexible framework that contains both the discrete and measure-theoretic inequalities treated here. 

These questions suggest that the inequalities studied in this paper can be viewed as a part of an even broader theory of weighted Hardy operators, nested averaging operators, and triangular kernel operators. The results obtained here thus give a foundation for that theory by separating the boundedness problem, which admits a complete characterization, from the sharper constant problem, for which local and structural methods appear to still be essential.




%%%%%%%%%%%%%%%%%%%%%%%%%%%%%%%%%%%%%%%%%%%%%%
%%%%%%%%%%%%% OLD REMARK %%%%%%%%%%%%%%%%%%%%%
%%%%%%%%%%%%%%%%%%%%%%%%%%%%%%%%%%%%%%%%%%%%%%

% \begin{rem}
%     If $m_n=1$ for all $n\geq 1$, then $w_n = |N_n|^{1/q}$. Moreover, if we know that $|N_n|$ is an increasing sequence, then we actually have
%     \begin{align*}
%         \rho = \sup_{n \geq 1}  \left(\frac{w_n}{M_{n}\!\left(1-\frac{w_n}{p w_{n+1}}\right)-\frac{M_{n-1}}{q}}\right) &= \sup_{n \geq 1}  \left(\frac{|N_{n}|^{1/q}}{M_{n}\!\left(1-\frac{|N_{n}|^{1/q}}{p\left|N_{n+1}\right|^{1/q}}\right)-\frac{M_{n-1}}{q}}\right) \\
%         %
%         &\leq \sup_{n \geq 1}  \left(\frac{|N_{n}|^{1/q}}{M_{n}\!\left(1-\frac{1}{p}\right)-\frac{M_{n-1}}{q}}\right) \\
%         %
%         &= q\cdot \sup_{n \geq 1}  \left(\frac{|N_{n}|^{1/q}}{M_{n}-M_{n-1}}\right),
%     \end{align*}
%     which is quite simple and elegant, although a bit less general and not sharp anymore in the case of \Cref{example 2}. 
% %
%     Lastly, observe that in the natural choice where $M_n:=|\mathbf{N}_n| = \sum_{k=1}^n |N_n|$, we actually obtain
%     \begin{align*}
%         \rho \leq q\cdot \sup_{n \geq 1}  \left(\frac{|N_{n}|^{1/q}}{M_{n}-M_{n-1}}\right)  =  q\cdot \sup_{n \geq 1}  \left(\frac{|N_{n}|^{1/q}}{|N_{n}|}\right) = \frac{q}{\inf_{n\geq 1}|N_n|^{1/p}} = \frac{q}{|N_1|^{1/p}}
%     \end{align*} 
% \end{rem}







% {\color{red} 
% \bigskip
% \textbf{QUESTION:} (1) Is the inequality optimal in other cases? How does it compare to previously known results? 

% \bigskip
% \textbf{QUESTION:} (2) In our previous paper, we were only able to show that \Cref{example 2} was sharp in the case of $p=2$. With this new technique, can we prove the optimality for all $1< p <\infty$?

% \bigskip
% \textbf{QUESTION:} (3) Can we show that the new condition is stronger than the old (not in the sense that it is always smaller than the other, since we know it is not the case, but in the sense that if one diverges, than the other also diverges but not the other way around)?

% \bigskip
% (3) Can we add variable weights in two places, namely
% \begin{align*}
%     \sum_{n=1}^{\infty} \Bigg|\frac{1}{M_n}\sum_{k \in \mathbf{N}_n} \bm{m_k} a_k \Bigg|^p
%     \leq~ \rho^p  \sum_{n=1}^{\infty} \bm{w_n} |a_n|^p.
% \end{align*}

% \bigskip
% (4) Can we generalize to general function spaces?
% }

% \bigskip 
% (5) What about kernels. Can we find a constant $C > 0$ such that
%     $$
%         \int_X \left|\int_X k (x, y) f (y) \, \mathrm{d}\mu (y) \right|^p \, \mathrm{d}\mu (x) \leq C \int_X |f (x)|^p \, \mathrm{d}\mu (x) .
%     $$
% Maybe some cases are covered by Schur's test. 

% If we let $X = (0, \infty )$, $\mu = dx$, and
%     $$
%         k (x, y) := \begin{cases}
%             \frac{1}{x} & y \leq x \\ 
%             0 & y > x
%         \end{cases}
%     $$
% then for any $x > 0$, 
%     $$
%         \int_X k (x, y) f (y) \, \mathrm{d}\mu (y) = \frac{1}{x} \int_0^x f (y) \, dy
%     $$
% I don't think that we can apply Schur's test to obtain Hardy's inequality. 

% Let $X = \mathbb{N}$ and $\mu$ as the counting measure on $\mathbb{N}$. Let $\mathbb{N} := N_1 \cup N_2 \cup \cdots$ be a partition and let $\mathbf{N}_n := N_1 \cup \cdots \cup N_n$. Define the following kernel function:
%     $$
%         k (n, j) := \begin{cases}
%             \frac{1}{M_n} m_j & j \in \mathbf{N}_n \\ 
%             0 & j \not\in \mathbf{N}_n ,
%         \end{cases}
%     $$
% where $M_n$ and $m_j$ are weights as defined previously. Then for a sequence $f : \mathbb{N} \rightarrow \mathbb{C}$, we compute 
%     $$
%         \int_{\mathbb{N}} k (n, j) f(j) \, \mathrm{d}\mu (j) =  \frac{1}{M_n} \sum_{j \in \mathbf{N}_n} m_j f (j) .
%     $$
% A special case we can try to prove is the following inequality:
%     $$
%         \int_X \left| \frac{1}{M (x)} \int_X m (y) f (y) \, \mathrm{d}\mu (y) \right|^p \, \mathrm{d}\mu (x) \leq C^p \int_X |f (x)|^p \, \mathrm{d}\mu (x) ,
%     $$
% where $M : X \rightarrow (0, \infty )$, $m : X \rightarrow (0, \infty )$ are two positive measurable functions, and 
%     $$ 
%         C :=
%     $$

% \bigskip 
% (6) Generalize to a sum over some integrals. There are two ways of doing that (probably). 

%     \begin{enumerate}[label=\textbullet]
%         \item {\color{blue} Let $X$ be a measure space with measure $\mu$ (ignore for now some technical assumptions on $X$). Let $(b_n)$ and $(M_n)$ be sequences of positive numbers and let $m$ be a positive measurable function defined on $X$. Let $X = X_1 \cup X_2 \cup \cdots$ be a partition of $X$ and set $X^{(n)} := X_1 \cup \cdots \cup X_n$. Let also
%             $$
%                 w_n := \Big(\int_{X_n} |m (x)|^q \, \mathrm{d}\mu (x) \Big)^{1/q}
%             $$
%         Then the inequality would look like
%             $$
%                 \sum_{n = 1} b_n \left| \frac{1}{M_n} \int_{\mathbf{N}_n} m (x) f (x) \, \mathrm{d}\mu (x) \right|^p \leq \rho^p \sum_{n = 1}^\infty b_n \int_{X_n} |f (x)|^p \, \mathrm{d}\mu (x) 
%             $$
%         where
%             $$
%                 \rho := \inf_{0 < t < \infty} \sup_{n \geq 1} \frac{b_n}{\frac{b_n M_n}{w_n} - \frac{b_{n + 1} M_n}{p w_{n + 1} t^q} - \frac{p-1}{p} \frac{b_n M_{n-1} t^p}{w_n}}
%             $$
%         I think the proof you have can be generalized to that case.}
%         \item Another approach would be to take, instead of one measure space, some finite measure spaces $(X_n, \mu_n )$. Let $(b_n)_{n\geq 1}$, $(M_n)_{n\geq1}$ be as before, let $m_n$ be positive bounded measurable functions on $X_n$ and let
%             $$
%                 w_n := \Big(\int_{X_n} |m_n (x)|^q \, \mathrm{d}\mu_n (x) \Big)^{1/q} .
%             $$
%         Then the inequality in this case would look like
%             $$
%                 \sum_{n = 1} b_n \left| \frac{1}{M_n} \sum_{k = 1}^n \int_{X_k} m_k (x) f (x) \, \mathrm{d}\mu_k (x)  \right|^p \leq \rho^p \sum_{n = 1}^\infty b_n \int_{X_n} |f (x)|^p \, \mathrm{d}\mu_n (x) 
%             $$
%         with the same $\rho$. We might need to be careful when applying some of the arguments used in your proof of the discrete case because we are in different measure spaces. \textcolor{red}{[I think this would work, but the case where all $b_n$ are equal to $1$ (which is the main case we are interested in) no longer gives something nice on the right-hand side.]} 
%     \end{enumerate}





















% \newpage

% \begin{problem}
%     Do there exist sequences $\alpha, \beta, \gamma \in \ell^{2}(\mathbb{Z}^{+})$, not identically zero, such that: 
%     \begin{enumerate} 
%     \item $\displaystyle \sum_{j = k}^{\infty} \overline{\gamma}_{j-k}\,\alpha_{j} = \beta_{k} \qquad \forall k \in \mathbb{Z}^{+}, $ 
%     \item $\displaystyle \sum_{k = 0}^{j} \gamma_{j-k}\,\beta_{k} = -\alpha_{j} \qquad \forall j \in \mathbb{Z}^{+} \text{ with }\alpha_{j} \neq 0? $
%     \end{enumerate}
% \end{problem}

% \noindent\textbf{Answer.} No.

% \begin{proof}
% Suppose that $\alpha,\beta,\gamma\in \ell^{2}(\mathbb{Z}^{+})$ satisfy the two identities.  
% Take the complex conjugate of (1):
% \[
% \overline{\beta_k}
%     = \sum_{j=k}^{\infty} \gamma_{ j-k}\,\overline{\alpha_j},
%     \qquad \forall k\ge 0 .
% \]
% Consider the inner product
% \[
% \langle \gamma * \beta + \alpha ,\, \alpha\rangle
%     = \sum_{j=0}^{\infty} \big( (\gamma * \beta)_j + \alpha_j \big)\,\overline{\alpha_j},
% \]
% where $(\gamma*\beta)_j := \sum_{k=0}^j \gamma_{\, j-k} \beta_k$. 
% %
% On the indices where $\alpha_j\neq 0$, property (2) gives 
% \[
% (\gamma * \beta)_j = -\alpha_j,
% \]
% and when $\alpha_j=0$ the summand in the inner product is zero.  
% Thus
% \[
% \langle \gamma * \beta + \alpha ,\, \alpha\rangle = 0.
% \tag{$*$}
% \]

% On the other hand,
% \[
% \langle \gamma * \beta ,\, \alpha\rangle
% = \sum_{j=0}^{\infty}\left( \sum_{k=0}^{j} \gamma_{j-k}\beta_k \right)\overline{\alpha_j} = \lim_{N\to\infty}\sum_{j=0}^{N}\left( \sum_{k=0}^{j} \gamma_{j-k}\beta_k \right)\overline{\alpha_j}.
% \]
% Hence, if we write $\beta^{(N)} := (\beta_0,\beta_1,\dots,\beta_N,0,0,\dots)$, then
% \begin{align*}
%     \lim_{N\to\infty}\sum_{j=0}^{N}\left( \sum_{k=0}^{j} \gamma_{j-k}\beta_k \right)\overline{\alpha_j} &= \sum_{k=0}^{N} \beta_k \!\left( \sum_{j=k}^{N} \gamma_{j-k}\, \overline{\alpha_j} \right)
% \end{align*}



% \[
% \begin{aligned}
% \langle \gamma * \beta ,\, \alpha\rangle
% &= \sum_{j=0}^{\infty}\left( \sum_{k=0}^{j} \gamma_{j-k}\beta_k \right)\overline{\alpha_j} \\
% &= \sum_{k=0}^{\infty} \beta_k \!\left( \sum_{j=k}^{\infty} \gamma_{j-k}\, \overline{\alpha_j} \right) \\
% &= \sum_{k=0}^{\infty} \beta_k\,\overline{\beta_k}
%     = \|\beta\|_2^2,
% \end{aligned}
% \]
% where we used the conjugated form of (1) in the second step.  
% Therefore,
% \[
% \langle \gamma * \beta + \alpha ,\, \alpha\rangle
% = \langle \gamma * \beta,\, \alpha\rangle + \|\alpha\|_2^2
% = \|\beta\|_2^2 + \|\alpha\|_2^2.
% \tag{$\dagger$}
% \]

% Equating ($*$) and ($\dagger$) gives
% \[
% \|\beta\|_2^2 + \|\alpha\|_2^2 = 0,
% \]
% so $\alpha = \beta = 0$.  %For these values, any $\gamma\in\ell^2(\mathbb{Z^+)}$ will satisfy (1) and (2).  

% Hence the only solution in $\ell^{2}(\mathbb{Z}^{+})$ is the trivial one.
% \end{proof}


% Define the operator $T_\gamma: \ell^2 \to \ell^2$ by
% \[
% (T_\gamma \alpha)_k := \sum_{j=k}^{\infty} \overline{\gamma}_{j-k}\, \alpha_j, \quad k \in \mathbb{Z}^+.
% \]
% Then the equation (1) becomes
% \[
% \beta = T_\gamma \alpha.
% \]
% Moreover, the condition (2) can be interpreted as
% \[
% \alpha = - T_\gamma^* \beta,
% \]
% where $T_\gamma^*$ is the adjoint of $T_\gamma$ on $\ell^2$, given by
% \[
% (T_\gamma^* \beta)_j = \sum_{k=0}^{j} \gamma_{j-k}\, \beta_k.
% \]
% Composing the two equations gives
% \[
% \alpha = - T_\gamma^* \beta = - T_\gamma^* T_\gamma \alpha,
% \]
% or equivalently,
% \[
% (I + T_\gamma^* T_\gamma) \alpha = 0.
% \]





% \begin{proof}
% Let $H^{2}$ be the Hardy space on the unit disc $\mathbb{D}$, identified with
% $\ell^{2}(\mathbb{Z}^{+})$ by the coefficient map
% \[
% A(z)=\sum_{j\ge 0}\alpha_j z^j
% \quad\longleftrightarrow\quad
% \alpha=(\alpha_j)_{j\ge 0}\in\ell^{2}(\mathbb{Z}^{+}),
% \]
% and similarly for $B(z)=\sum_{k\ge 0}\beta_k z^k$ and $G(z)=\sum_{n\ge 0}\gamma_n z^n$.

% Assume $\gamma\in\ell^{2}(\mathbb{Z}^{+})$ is not identically zero, hence
% $G\in H^{2}$.
% Since $G$ need not belong to $H^{\infty}$, the Toeplitz operators below may be
% unbounded; we work on their canonical dense domains.

% Define the (possibly unbounded) Toeplitz operators with analytic and co-analytic
% symbols by
% \[
% T_{G}f := P(Gf), 
% \qquad 
% T_{\overline{G}}f := P(\overline{G}f),
% \]
% with domains
% \[
% \mathcal{D}(T_{G})
% :=\{f\in H^{2}: Gf\in L^{2}(\mathbb{T})\},
% \qquad
% \mathcal{D}(T_{\overline{G}})
% :=\{f\in H^{2}: \overline{G}f\in L^{2}(\mathbb{T})\},
% \]
% where $P:L^{2}(\mathbb{T})\to H^{2}$ is the Szeg\H{o} projection.

% Suppose there exists $A\in \mathcal{D}(T_{\overline{G}})$, $A\neq 0$, such that
% \begin{equation}\label{eq:kernel}
% A\in\mathcal{D}(T_{G}T_{\overline{G}})
% \quad\text{and}\quad
% T_{G}T_{\overline{G}}A=-A .
% \end{equation}
% Set
% \begin{equation}\label{eq:Bdef}
% B:=T_{\overline{G}}A .
% \end{equation}
% We prove that the coefficient sequences
% $\alpha,\beta,\gamma$ of $A,B,G$ satisfy (1) and (2).

% \smallskip
% \noindent\textbf{Step 1: (1) is exactly $B=T_{\overline{G}}A$.}
% Write
% \[
% \overline{G(e^{it})}A(e^{it})
% =
% \Big(\sum_{n\ge 0}\overline{\gamma_n}e^{-int}\Big)
% \Big(\sum_{j\ge 0}\alpha_j e^{ijt}\Big)
% =
% \sum_{k\in\mathbb{Z}}
% \Big(\sum_{j-k\ge 0}\overline{\gamma_{j-k}}\alpha_j\Big)e^{ikt}.
% \]
% Projecting onto nonnegative frequencies gives
% \[
% B(e^{it})=P(\overline{G}A)(e^{it})
% =\sum_{k\ge 0}
% \Big(\sum_{j=k}^{\infty}\overline{\gamma_{j-k}}\alpha_j\Big)e^{ikt}.
% \]
% Since $B(z)=\sum_{k\ge 0}\beta_k z^k$, we obtain
% \[
% \beta_k=\sum_{j=k}^{\infty}\overline{\gamma_{j-k}}\alpha_j
% \qquad\forall k\ge 0,
% \]
% which is exactly (1).

% \smallskip
% \noindent\textbf{Step 2: (2) is exactly $T_{G}B=-A$ on the support of $\alpha$.}
% From \eqref{eq:kernel} and \eqref{eq:Bdef},
% \[
% T_{G}B=T_{G}T_{\overline{G}}A=-A .
% \]
% Expand $GB$ in Fourier series:
% \[
% G(e^{it})B(e^{it})
% =
% \Big(\sum_{m\ge 0}\gamma_m e^{imt}\Big)
% \Big(\sum_{k\ge 0}\beta_k e^{ikt}\Big)
% =
% \sum_{j\ge 0}
% \Big(\sum_{k=0}^{j}\gamma_{j-k}\beta_k\Big)e^{ijt}
% +\ (\text{negative frequencies}).
% \]
% Projecting onto $H^{2}$ gives
% \[
% T_{G}B(e^{it})
% =
% \sum_{j\ge 0}
% \Big(\sum_{k=0}^{j}\gamma_{j-k}\beta_k\Big)e^{ijt}.
% \]
% But $T_{G}B=-A=\sum_{j\ge 0}(-\alpha_j)e^{ijt}$, hence
% \[
% \sum_{k=0}^{j}\gamma_{j-k}\beta_k=-\alpha_j
% \qquad\forall j\ge 0.
% \]
% In particular this holds for every $j$ with $\alpha_j\neq 0$, which is (2).

% \smallskip
% Thus, any nonzero $A$ satisfying \eqref{eq:kernel} yields
% nontrivial sequences $\alpha,\beta,\gamma\in\ell^{2}(\mathbb{Z}^{+})$
% satisfying (1) and (2).
% \end{proof}





% Using (2) we have
% \[
% \sum_{k = 0}^{j} \gamma_{j-k}\,\beta_{k} = -\alpha_{j}
% \]
% for all $\alpha_j\neq 0$. Hence, multiplying both sides by $\overline{\alpha_j}$, we have
% \[
% \sum_{k = 0}^{j} \gamma_{j-k}\,\beta_{k} \overline{\alpha_j} = -|\alpha_{j}|^2.
% \]
% Summing with respect to $j$ yield
% \[
% \sum_{j=0}^\infty \sum_{k = 0}^{j} \gamma_{j-k}\,\beta_{k} \overline{\alpha_j} = -\sum_{j=0}^\infty|\alpha_{j}|^2 = -\|\alpha\|^2.
% \]


% Since $\gamma\in \ell^2(\mathbb{Z)}$, each entry in the sequence is bounded, say by $M>0$. Hence, we have
% \begin{align*}
%     \sum_{j=0}^{\infty} \sum_{k=0}^{j} |\gamma_{j-k}\beta_k \overline{\alpha_j}| \leq \sum_{j=0}^{\infty} |\sum_{k=0}^{j} \gamma_{j-k}\beta_k \overline{\alpha_j}|
% \end{align*}


    

    \bibliographystyle{plain}
    \bibliography{ref}

\end{document}


